\documentclass[11pt]{report}

% ---------------------------------------------------------------------------
% normet User Manual
% Build:  docs/manual/build.sh   (runs pdflatex until the references settle)
% Figures: docs/manual/make_figures.py   (screenshots of the real window)
% ---------------------------------------------------------------------------
\usepackage[utf8]{inputenc}
\usepackage[T1]{fontenc}
\usepackage{lmodern}
% Typewriter text is code: "--" stays two hyphens, as a command-line option is typed,
% instead of becoming an en dash.
\usepackage{microtype}
\DisableLigatures[-]{encoding = T1, family = tt*}
\usepackage[a4paper,margin=1in]{geometry}
\usepackage{amsmath,amssymb}
\usepackage{graphicx}
\usepackage{booktabs}
\usepackage{longtable}
\usepackage{array}
\usepackage{xcolor}
\usepackage{fancyhdr}
\usepackage{textcomp}
\usepackage{listings}
\usepackage[hidelinks,colorlinks=true,linkcolor=blue!55!black,urlcolor=blue!55!black,
            citecolor=blue!55!black]{hyperref}

% Compact list spacing without the enumitem package.
\newcommand{\nosep}{\setlength{\itemsep}{0pt}\setlength{\parskip}{0pt}\setlength{\topsep}{2pt}}

\definecolor{nmblue}{HTML}{2C7BB6}
\definecolor{nmgray}{HTML}{8E8E93}
\definecolor{codebg}{HTML}{F5F5F7}
\definecolor{warn}{HTML}{FF9500}

% --- code listings -----------------------------------------------------------
\lstdefinestyle{normet}{
  backgroundcolor=\color{codebg},
  basicstyle=\ttfamily\small,
  keywordstyle=\color{nmblue}\bfseries,
  commentstyle=\color{nmgray}\itshape,
  stringstyle=\color{green!45!black},
  breaklines=true, frame=single, rulecolor=\color{nmgray!40},
  showstringspaces=false, upquote=true, columns=fullflexible, xleftmargin=4pt, framexleftmargin=4pt,
  literate={−}{{$-$}}1 {≥}{{$\geq$}}1 {→}{{$\rightarrow$}}1,
}
\lstset{style=normet}

% --- headers / footers -------------------------------------------------------
\pagestyle{fancy}
\fancyhf{}
\fancyhead[L]{\small\textsc{normet User Manual}}
\fancyhead[R]{\small\leftmark}
\fancyfoot[C]{\thepage}
\renewcommand{\headrulewidth}{0.4pt}

% --- handy macros ------------------------------------------------------------
\newcommand{\nm}{\textsf{normet}}
% Code may break after an underscore, where a break reads naturally in a name, so a long
% identifier near the end of a line wraps instead of running into the margin.
\newcommand{\codeunderscore}{\textunderscore\allowbreak}
\newcommand{\code}[1]{\texttt{\let\_\codeunderscore #1}}
% A line that still cannot be set is stretched a little rather than left overfull.
\setlength{\emergencystretch}{2.5em}
\newcommand{\tab}[1]{\textbf{#1}}
\newcommand{\NOtwo}{NO\textsubscript{2}}
\newcommand{\PMtwo}{PM\textsubscript{2.5}}
\newcommand{\ugm}{\ensuremath{\mu}g\,m\textsuperscript{$-3$}}
\newcommand{\note}[1]{\par\noindent\colorbox{codebg}{\parbox{0.97\linewidth}{\small\textbf{Note.} #1}}\par}
\newcommand{\caution}[1]{\par\noindent\colorbox{warn!15}{\parbox{0.97\linewidth}{\small\textbf{Caution.} #1}}\par}
\newcommand{\screenshot}[3]{%
  \begin{figure}[htbp]\centering
    \fboxsep=0pt\fbox{\includegraphics[width=0.98\linewidth]{figures/#1}}%
    \caption{#2}\label{#3}
  \end{figure}}

\begin{document}

% ============================ TITLE PAGE ====================================
\begin{titlepage}
\centering
\vspace*{2.5cm}
{\IfFileExists{../../packaging/assets/normet.png}{\includegraphics[width=3.2cm]{../../packaging/assets/normet.png}\\[1cm]}{}}
{\Huge\bfseries\color{nmblue} normet\par}
\vspace{0.4cm}
{\Large Normalisation, Decomposition and Counterfactual Modelling\\for Environmental Time Series\par}
\vspace{1.2cm}
{\large\bfseries User Manual\par}
\vspace{0.3cm}
{\large Version 1.0.0\par}
\vfill
{\large A desktop application and Python library for removing the weather\\
from air-quality records and asking what an intervention changed.\par}
\vspace{1.5cm}
{\color{nmgray}\today}
\end{titlepage}

\tableofcontents
\cleardoublepage

% ============================ 1. INTRODUCTION ===============================
\chapter{Introduction}\label{ch:intro}

\section{What is normet?}
\nm{} is a desktop application and Python library for the analysis of environmental
time series, air-quality records above all. It answers two questions that a raw
concentration record cannot answer by itself:
\begin{enumerate}\nosep
  \item \emph{How much of the change I see is the weather?} A month of stagnant,
        cold anticyclonic air raises concentrations as surely as a new source does.
        \nm{} learns how the pollutant depends on the meteorology and the calendar, then
        averages the meteorology out, leaving the part of the record the weather does
        not explain. This is \emph{meteorological normalisation}, or
        \emph{de-weathering}.
  \item \emph{What would have happened without the intervention?} A low-emission
        zone, a lockdown or a plant closure changes one place at one date. \nm{}
        builds a counterfactual for that place from other places the intervention did
        not touch --- a \emph{synthetic control} --- and measures the difference after
        the date, with placebo tests and uncertainty bands to say whether the
        difference is larger than chance.
\end{enumerate}
Everything the window does is a call into the \nm{} Python package, which can be
scripted, run from the command line, or driven by an AI assistant over the Model
Context Protocol (Chapter~\ref{ch:script}). The window writes down each call it makes,
so an analysis done by clicking can be saved as a script that repeats it
(Section~\ref{sec:commandlog}).

\section{The idea in one equation}
Let $y_t$ be the concentration at time $t$, $\mathbf{m}_t$ the meteorology at that time
(wind speed and direction, temperature, humidity, boundary-layer height, \dots) and
$\mathbf{c}_t$ the calendar (a trend counter, the day of the year, the weekday, the hour).
\nm{} fits a machine-learning model
\begin{equation}
  y_t \approx f(\mathbf{m}_t, \mathbf{c}_t)
\end{equation}
and then asks what $f$ predicts at time $t$ if the weather at $t$ had been any other
weather the site has seen. Averaging over many such draws,
\begin{equation}
  \tilde{y}_t \;=\; \frac{1}{N}\sum_{k=1}^{N} f\bigl(\mathbf{m}^{(k)}, \mathbf{c}_t\bigr),
  \qquad \mathbf{m}^{(k)} \sim \text{the record's own meteorology},
  \label{eq:norm}
\end{equation}
gives the \emph{normalised} series $\tilde{y}_t$: the concentration expected at time $t$
under typical weather. Because the calendar $\mathbf{c}_t$ is kept, a change in
emissions --- which the model can only express through the trend, season, week and
hour --- survives in $\tilde{y}_t$, while the day-to-day swings that the weather caused
are averaged away. Chapter~\ref{ch:theory} develops this and the decompositions built on
it.

\section{Key features}
\begin{itemize}\nosep
  \item AutoML training with FLAML (an estimator search within a time budget) or a
        tuned LightGBM, with the held-out fit reported as $R^2$, RMSE and nine
        further statistics.
  \item Monte-Carlo meteorological normalisation, with a choice of which variables
        are resampled, uncertainty quantiles, reference pools and conditional
        resampling; an adaptive version that stops when the series has converged.
  \item Emission decomposition (trend, seasonal, weekly, diurnal, short-term) and
        meteorological decomposition (one contribution per variable or named group,
        sequential or Shapley).
  \item Rolling (windowed) normalisation, and a multi-scale decomposition that
        separates the weather's influence by timescale.
  \item Partial-dependence curves of the trained model.
  \item Zero-shot de-weathering with the Chronos-2 time-series foundation model, on
        the same result schema --- no training needed.
  \item Synthetic control on panel data: augmented SCM, classic Abadie SCM,
        difference-in-differences, matrix completion, robust SCM, ML-SCM and Bayesian
        SCM, with placebo-in-space and placebo-in-time tests, jackknife and bootstrap
        bands, and an all-units batch.
  \item Data Studio: air quality from the six UK networks, the EEA reporting network
        for Europe, or OpenAQ for most of the world, with reanalysis meteorology,
        merged into a model-ready table (no key for the UK networks or the EEA).
  \item Transport Studio: HYSPLIT back-trajectories turned into transport predictors
        and joined onto the record.
  \item A command log that writes every action as the call that repeats it; HTML
        reports with provenance; figures laid out for a journal's column.
  \item An Assistant (a language model working in the window) that proposes the
        target and predictors and reads the results with you.
\end{itemize}

\section{How this manual is organised}
Chapter~\ref{ch:install} covers installation. Chapter~\ref{ch:theory} explains the
methods in enough depth to interpret every number the program reports. Chapters
\ref{ch:start}--\ref{ch:pdp} walk through a complete de-weathering analysis in the
window on a real record, and Chapter~\ref{ch:scm} the synthetic-control workflow on a
real panel. Chapter~\ref{ch:data} covers getting data in, Chapter~\ref{ch:export}
exporting, reporting and reproducing a session, Chapter~\ref{ch:assistant} the
Assistant and Chapter~\ref{ch:script} scripting. The appendices list the defaults, where
\nm{} keeps its files, and a troubleshooting guide.

The figures are screenshots of the real window, made by
\code{docs/manual/make\_figures.py} on the record \nm{} ships with: hourly \NOtwo{} and
ERA5 meteorology at London Marylebone Road (site MY1), January to August 2020, a window
that spans the UK's first COVID-19 lockdown. The numbers quoted in the text are the ones
that script printed. A run of your own will differ in the last digit or two: the AutoML
search is bounded by wall-clock time, so it does not visit exactly the same candidates
twice.

This manual is also published in Chinese (\code{normet\_manual\_zh.pdf}), chapter for
chapter, with its screenshots of the window in Chinese (Section~\ref{sec:appearance}).

% ============================ 2. INSTALLATION ===============================
\chapter{Installation}\label{ch:install}

\nm{} ships as a self-contained desktop application for each platform, and as a Python
package on PyPI. Choose the installers for interactive analysis, the package for
scripting, servers and notebooks.

\section{macOS}
Open \code{Normet-1.0.0-macos-arm64.dmg} and drag \code{Normet.app} onto the
\code{Applications} folder. The application carries its own Python, Qt, FLAML, LightGBM
and XGBoost, and the Data Studio's readers, so nothing else is needed for the
de-weathering, synthetic-control and data workflows. The first launch takes longer than
the rest while it builds a font cache, which later launches reuse (it is kept in
\code{\textasciitilde/.normet/matplotlib}); a splash screen shows meanwhile.

\note{The application is not signed with an Apple Developer ID. The first time, open it
by right-clicking \code{Normet.app} and choosing \emph{Open}, then confirming in the
dialog; this is needed only once. Alternatively run
\code{xattr -dr com.apple.quarantine /Applications/Normet.app} in Terminal.}

The installers do not include the optional Chronos-2 foundation model (it needs PyTorch,
about a gigabyte on its own), the Bayesian SCM backend (PyMC) or the MCP server. Their
options are greyed out in the window, with a tooltip naming what is missing; install
\nm{} from PyPI with the extra to use them (below).

\section{Windows and Linux}
On Windows run \code{normet-setup-1.0.0.exe}, an Inno~Setup wizard that adds Start-menu
and optional desktop shortcuts. The installer is unsigned, so SmartScreen shows
\emph{Windows protected your PC} on first run: choose \emph{More info} $\rightarrow$
\emph{Run anyway}. On Linux download \code{normet-1.0.0-x86\_64.AppImage}, mark it
executable (\code{chmod +x}) and run it. On a minimal distribution Qt may need
\code{libegl1 libgl1 libxkbcommon0 libdbus-1-3 libxcb-cursor0}.

The installers are built by the repository's \code{build-gui.yml} workflow (one per
operating system; PyInstaller cannot cross-compile) and are attached to its runs under
\emph{Actions $\rightarrow$ Artifacts}. \code{packaging/README.md} explains how to build
one locally.

\section{From PyPI}
\nm{} needs Python~3.10 or newer (3.10--3.12 are tested) and NumPy below 2. The core
package is small; everything heavier is an extra:
\begin{lstlisting}[language=bash]
python -m venv .venv && source .venv/bin/activate
pip install "normet[gui,flaml]"      # the desktop app with the AutoML backend
normet-gui                           # or: normet-gui mydata.csv
\end{lstlisting}

\begin{center}\small
\begin{tabular}{@{}ll@{}}
\toprule
Extra & What it adds \\
\midrule
\code{flaml}      & FLAML AutoML backend (with XGBoost, which FLAML imports) --- recommended \\
\code{lgb}        & LightGBM backend with random hyperparameter search \\
\code{gui}        & PySide6, for \code{normet-gui} \\
\code{data}       & ERA5 (CDS), UK networks, OpenAQ, EEA, Open-Meteo adapters \\
\code{geo}        & Shapely and pyshp: source regions from GeoJSON/Shapefile \\
\code{cli}        & Click and PyYAML, for the \code{normet} command \\
\code{foundation} & Chronos-2 and PyTorch: zero-shot de-weathering, station embeddings \\
\code{finetune}   & \code{foundation} plus \code{accelerate} and \code{peft}, to fine-tune Chronos-2 \\
\code{bayes}      & PyMC and ArviZ, for the Bayesian SCM \\
\code{mcp}        & The MCP server \code{normet-mcp}, for AI assistants \\
\code{dask}       & Dask, for large data sets \\
\code{all}        & Everything above \\
\bottomrule
\end{tabular}
\end{center}

Importing \nm{} never requires an optional dependency: a function that needs one says
which extra to install when it is called.

% ============================ 3. METHODOLOGY =================================
\chapter{Methodology}\label{ch:theory}

This chapter explains what \nm{} computes, so that the numbers it reports can be
interpreted with confidence. You do not need it to run an analysis; you do need it to
defend one.

\section{Preparing the data}\label{sec:prepare}
Before a model is trained the table goes through \code{prepare\_data}, in this order:
\begin{enumerate}\nosep
  \item \textbf{A date column.} The table must have one datetime column called
        \code{date} (or an index, or a column named \code{datetime}, \code{time} or
        \code{timestamp}, which is renamed). Dates with time zones are converted to UTC
        and then made naive.
  \item \textbf{Target and predictors.} The target is renamed \code{value}; the table is
        reduced to the date, the target and the predictors.
  \item \textbf{Missing values.} Rows with a missing target are dropped. A missing
        numeric predictor is filled with that column's median; a missing categorical
        one with its most frequent value.
  \item \textbf{Calendar variables.} Four columns are added: \code{date\_unix} (seconds
        since 1970, the trend), \code{day\_julian} (day of the year, the season),
        \code{weekday} (1~=~Monday to 7~=~Sunday, categorical) and \code{hour}.
        The model uses only those listed among its predictors; the window always lists
        all four.
  \item \textbf{Train/test split.} A column \code{set} marks each row
        \code{training} or \code{testing} (Section~\ref{sec:split}).
\end{enumerate}

\subsection{The split}\label{sec:split}
The held-out rows are what the model's fit is judged on. Four splits are offered:
\begin{description}\nosep
  \item[\code{random}] (default) A random 25\,\% of rows are held out. Every season,
        weekday and hour is represented in training.
  \item[\code{ts}] The last 25\,\% of the record is held out. A trend variable must then
        be extrapolated past anything seen in training, which tree models cannot do;
        use it to test forecasting, not to train for normalisation.
  \item[\code{month\_ts}] Within every calendar month of every year, one contiguous
        block of 25\,\% is held out, at a random (seeded) position in the month.
  \item[\code{season\_ts}] The same within every meteorological season (DJF, MAM, JJA,
        SON; December counts with the following year's winter).
\end{description}

\caution{Under \code{random}, neighbouring hours fall on both sides of the split, so an
autocorrelated record looks easier to predict than it is: the test $R^2$ is an upper
bound. This matters most for predictors that are constant over several hours, such as
back-trajectory features released every six hours, which a tree model can use as a time
stamp. On the bundled MY1 \PMtwo{} case a trajectory-only model scores about 0.84 under a
random split and 0.46 under a blocked one. Judge such models with \code{month\_ts}.}

\section{Training the model}\label{sec:train}
\nm{} trains gradient-boosted trees (by default) that map the predictors to the target.
Two backends are offered:
\begin{description}\nosep
  \item[FLAML] An AutoML search over the estimators ticked (LightGBM alone by default;
        XGBoost, depth-limited XGBoost, random forest, extra trees and scikit-learn's
        histogram gradient boosting can be added) within a wall-clock \emph{time budget},
        maximising $R^2$. The library's default budget is 90~s; the window's is 60~s.
  \item[LightGBM] LightGBM with a random hyperparameter search of a set number of
        \emph{trials} (20 in the window).
\end{description}
A longer budget rarely hurts and often helps on long records. Because the FLAML budget is
wall-clock time, the same script on the same data can stop at a different candidate on a
busier machine; set \code{max\_iter} in \code{model\_config} with \code{time\_budget=-1}
for a search bounded by trials instead.

\subsection{Model statistics}\label{sec:modstats}
\code{modStats} predicts the training rows, the testing rows and all rows, and reports
for each: $n$; FAC2 (the fraction of predictions within a factor of two of the
observation); MB and MGE (mean bias and mean gross error); NMB and NMGE (the same,
normalised by the sum of observations); RMSE; $r$ and its significance; COE (coefficient
of efficiency); IOA (index of agreement); and $R^2$. The held-out ($\text{set} =
\text{testing}$) row is the one to judge the model by.

\section{Meteorological normalisation}\label{sec:normalise}
Equation~\eqref{eq:norm} is computed by Monte Carlo. For each of $N$ draws
(\code{n\_samples}, 300 by default):
\begin{enumerate}\nosep
  \item The \emph{resampled variables} of every row are replaced, all together, by those
        of a row drawn at random (with replacement) from the \emph{resample pool} ---
        by default the record itself. Whole rows are drawn, so the resampled variables
        keep their joint structure: a draw never pairs a winter temperature with a
        summer boundary layer.
  \item Every other predictor --- by default the four calendar variables --- keeps its
        observed value.
  \item The model predicts every row.
\end{enumerate}
The normalised series is the mean of the $N$ predictions at each timestamp; with
\code{return\_quantiles} the chosen quantiles of the $N$ predictions are returned as
well, as columns \code{q050}, \code{q950} and so on (the quantile $\times 1000$). These
quantiles describe the spread of the prediction over weather, not the uncertainty of the
model.

\paragraph{Which variables to resample.} Resampling the meteorology and keeping the
calendar is the usual de-weathering. Resampling a calendar variable too removes what it
carries: resampling \code{day\_julian} removes the seasonal cycle, resampling
\code{date\_unix} removes the trend. A predictor that is not meteorological --- traffic
counts, say --- can be kept at its observed values by leaving it unticked.

\paragraph{How many samples.} The Monte-Carlo error of each timestamp shrinks as
$1/\sqrt{N}$. 300 is ample for a smooth daily or monthly mean; individual hours are
noisier. \code{normalise\_auto} chooses $N$ for you: it adds batches of 100 draws until,
for three checks in a row, 95\,\% of the timestamps have a Monte-Carlo relative standard
error below 3\,\%, or 5000 draws have been made (a warning says when the cap was hit).

\paragraph{Reference pools.} \code{resample\_df} draws the weather from another table ---
2019's weather, for example, to ask what 2020 would have looked like under 2019's
meteorology. \code{conditional\_on} restricts the pool to rows meeting a condition --- a value, a list
of values, or a function of the column, as \code{\{"ws": lambda v: v < 2\}} draws from
calm hours only. \code{resample\_pools} draws some variables from a pool of
their own: \code{\{"transport": clean\_hours[traj\_cols]\}} draws the trajectory features
from clean-air-mass hours while the local weather is still drawn from the whole record.
Each pool has its own random stream, so adding or renaming a pool does not move the draws
of the others.

\paragraph{Reproducibility.} The draws are seeded (\code{seed=7654321} by default), so the
same model, data and settings give the same series to the last digit.

\section{Emission decomposition}\label{sec:emission}
The emission decomposition splits the normalised series into the timescales the calendar
variables carry, by \emph{nested} normalisation. Five normalisations are run with the
same seed, fixing one more calendar variable each time, in this fixed order:
\begin{center}\small
\begin{tabular}{@{}llll@{}}
\toprule
Run & Held at observed values & Resampled & Gives \\
\midrule
$b$         & nothing                                  & everything & base \\
$u$         & \code{date\_unix}                        & the rest   & trend $= u - b$ \\
$j$         & \ldots{} + \code{day\_julian}            & the rest   & seasonal $= j - u$ \\
$w$         & \ldots{} + \code{weekday}                & the rest   & weekly $= w - j$ \\
$h$         & \ldots{} + \code{hour}                   & meteorology & diurnal $= h - w$ \\
\bottomrule
\end{tabular}
\end{center}
The last run, $h$, is the ordinary normalised series, reported as \code{emi\_total}.
The base run is split into its mean, \code{emi\_base}, and its deviation from that mean,
\code{emi\_noise}; so
\begin{equation}
  \text{emi\_total} = \text{emi\_base} + \text{date\_unix} + \text{day\_julian}
     + \text{weekday} + \text{hour} + \text{emi\_noise}.
\end{equation}
A calendar variable that is not among the model's predictors gets no column.

\caution{The order decides what each component can hold. The trend is computed while the
day of the year, weekday and hour are still being averaged over, so it cannot carry a
signal that recurs at the same date every year --- only a genuine drift. The seasonal
component is computed with each row's own trend value already fixed, so it is not a
climatology: it stays specific to its year and can register a one-off event, such as a
single year's holiday dip or a lockdown, despite its label. Look for a recurring calendar
effect in \code{day\_julian}, not \code{date\_unix}.}

\section{Meteorological decomposition}\label{sec:metdecomp}
The meteorological decomposition splits the part of the prediction the weather accounts
for among the meteorological variables, or among named groups of them. The calendar is
held at its observed values throughout. Write $v(S)$ for the normalised series with the
variables in $S$ held at their observed values and the rest resampled; $v(\emptyset)$ is
\code{emi\_total}, and $v(\text{all})$ is the model's prediction with nothing resampled.
The difference $v(\text{all}) - v(\emptyset)$ is split in one of two ways:
\begin{description}
  \item[Sequential] Fix the variables one at a time and report each step's change:
        $\phi_i = v(\{1,\dots,i\}) - v(\{1,\dots,i-1\})$. This takes $k+1$
        normalisations for $k$ variables, but each contribution is conditional on those
        fixed before it, so the split depends on the order --- an explicit
        \code{variable\_order}, else the model's feature importance, which can reshuffle
        when the model is refitted.
  \item[Shapley] Average each variable's marginal effect over every order it could be
        fixed in:
        \begin{equation}
          \phi_i = \sum_{S \not\ni i} \frac{|S|!\,(k-|S|-1)!}{k!}\,
                   \bigl[v(S\cup\{i\}) - v(S)\bigr].
        \end{equation}
        No order is privileged. The exact values need all $2^k$ coalitions, allowed up
        to $k=10$; beyond that, group the variables or pass \code{n\_permutations} for
        an estimate from random orders (drawn in antithetic pairs).
\end{description}
Either way the contributions add up exactly to $v(\text{all}) - v(\emptyset)$, and every
normalisation uses the same seed, so the differences between them are paired (common
random numbers) rather than the difference of two noisy estimates. The result also
carries \code{met\_total} = observed $-$ \code{emi\_total}, its mean \code{met\_base},
and \code{met\_noise} = \code{met\_total} $-$ \code{met\_base} $-\sum_i\phi_i$, the part
of the observation the model does not reproduce.

\paragraph{Groups.} \code{groups=\{"local": met\_cols, "transport": traj\_cols\}} gives one
column per group; a group is fixed and resampled as a unit, so its column includes the
interactions among its members. Every meteorological predictor must be in exactly one
group. With groups the default attribution is Shapley; without them it is sequential.

\section{Rolling normalisation}\label{sec:rolling}
\code{rolling} normalises the record window by window: windows of \code{window\_days}
days (14 by default), started every \code{rolling\_every} days (7), each normalised with
its weather drawn \emph{from that window only}. A short window therefore removes only the
weather that varies within it --- synoptic variability for a fortnight --- and keeps the
slower meteorological variation that a full-record normalisation would remove. The result
has one column per window, counted from 1: \code{window\_01}, \code{window\_02}, \dots\
(padded to three digits past 99 windows, so the names sort in the windows' order);
\code{window\_columns(rolled)} lists them. The window shows their mean, with the spread
across overlapping windows.

\section{Multi-scale decomposition}\label{sec:multiscale}
Write $Y_W$ for the mean of the rolling-normalised series at window width $W$, and
$Y_\infty$ for the ordinary normalised series, whose pool is the whole record. As $W$
grows the window's pool approaches the record's, and $Y_W \to Y_\infty$. Differences
between widths isolate the weather influence specific to each band of timescales,
much as wavelet detail coefficients do:
\begin{align}
  D_\text{fast} &= Y_\text{fast} - Y_\text{meso} && \text{synoptic, sub-monthly} \\
  D_\text{meso} &= Y_\text{meso} - Y_\text{slow} && \text{intra-seasonal} \\
  D_\text{slow} &= Y_\text{slow} - Y_\infty      && \text{what even the slow window's pool leaves}
\end{align}
The widths default to 14, 90 and 365~days. Each scale's windows are stepped every quarter
of the width, and a scale with fewer than two windows on the record is skipped, with a
note saying so.

\section{Partial dependence}\label{sec:pdptheory}
A partial-dependence curve shows how the model's prediction responds to one predictor
with the others averaged over the data: for each of \code{grid\_points} values spanning
the 1st to 99th percentile of the variable, every training row's value of it is set to
that value and the predictions are averaged. \code{pdp\_mean} is that average and
\code{pdp\_std} the spread of the predictions about it. A curve describes the model, not
the atmosphere; correlated predictors (temperature and boundary-layer height, say) are
varied one at a time, into combinations the record may never contain.

\section{Zero-shot de-weathering with Chronos-2}\label{sec:chronos}
Chronos-2 is a pretrained time-series foundation model. \nm{} uses it without training:
the meteorology enters through the model's covariate channel, and the weather is
marginalised by resampling exactly as in Equation~\eqref{eq:norm}. It is useful where
there is too little history to train on, at a site that has just come online, or as a
second estimate from a model that never saw the target's own record. Four properties
decide whether it is the right tool:
\begin{itemize}\nosep
  \item \textbf{The first \code{context\_length} rows are not a result.} The model needs
        that much history (2048 rows by default) before it can project, so those rows
        are seeded with the observed values and the two curves coincide there by
        construction. Give it years, not months.
  \item \textbf{Every Monte-Carlo sample is a full forward pass}, so the default is 8
        samples, not 300. The model runs on CUDA, Apple Silicon (\code{mps}) or the CPU,
        in that order of preference; the CPU is slow.
  \item \textbf{Position is time.} Chronos-2 counts rows, so the record must be on a
        regular grid; \code{to\_regular\_index} rebuilds it, leaving gaps as missing.
  \item \textbf{Covariate sensitivity is the quality gate.} Nothing was fitted, so there
        is no held-out fit and no feature importance. Instead the future weather is
        shuffled and the shift in the forecast's median measured. Below about 1\,\% of
        the prediction the model is reading the target's own history and ignoring the
        weather, and a ``normalised'' series from it means nothing.
\end{itemize}
The meteorological decomposition has a zero-shot path. The emission decomposition, the
rolling normalisation and partial dependence do not: they are defined on a fitted
model's response to the calendar variables, which Chronos-2 reads off the target's own
history instead.

\section{Synthetic control}\label{sec:scmtheory}
A synthetic control estimates the effect of an intervention on one \emph{treated} unit
(a site, a city) from a panel of units observed before and after a \emph{cutoff} date.
Units the intervention did not touch --- the \emph{donor pool} --- are combined into a
synthetic version of the treated unit that tracks it before the cutoff; after the cutoff,
the synthetic series is the counterfactual, and
\begin{equation}
  \text{effect}_t = \text{observed}_t - \text{synthetic}_t .
\end{equation}
The average effect over the post-cutoff period is reported as the ATT (average treatment
effect on the treated). For an air-quality outcome, a de-weathered series is the natural
input: the weather is then not part of what the donors must reproduce.

\subsection{Estimators}
\begin{description}\nosep
  \item[\code{scm}] Augmented SCM (the default). At each date a ridge regression across
        the donors predicts the outcome from each unit's pre-cutoff trajectory; the
        synthetic value is the treated unit's ridge prediction plus a weighted sum of the
        donors' ridge residuals, the weights summing to one (non-negative unless allowed
        otherwise).
  \item[\code{abadie}] Classic SCM: non-negative donor weights summing to one, fitted to
        the treated unit's pre-cutoff path. Interpretable, but it cannot reach a treated
        unit that lies outside the range of its donors.
  \item[\code{did}] Difference-in-differences: a two-way fixed-effects baseline that
        assumes parallel trends.
  \item[\code{mcnnm}] Matrix completion with nuclear-norm minimisation (Athey et~al.,
        2021): the post-cutoff treated cells are treated as missing and filled from a
        low-rank fit of the panel.
  \item[\code{robust}] Robust SCM (Amjad, Shah and Shen, 2018): the donor matrix is
        de-noised by singular-value thresholding before the weights are fitted.
  \item[\code{mlscm}] ML-augmented SCM (\emph{experimental}): an AutoML model maps donor
        outcomes to the treated outcome in the pre-period. It lacks the causal
        guarantees of the others and warns so.
  \item[\code{bayesian}] Bayesian simplex SCM (PyMC): a Dirichlet prior on the weights,
        giving posterior credible bands directly.
\end{description}

\subsection{Diagnostics}\label{sec:scmdiag}
\code{scm\_diagnostics} reports the pre-cutoff fit (\code{pre\_rmse}, \code{pre\_mae},
\code{pre\_mape}, \code{pre\_r2}), the post-cutoff effect (\code{att}, \code{att\_cum},
\code{post\_rmse}) and how concentrated the weights are: the Herfindahl index
$\text{HHI} = \sum_j w_j^2$, the effective number of donors $1/\text{HHI}$, and the top
donors with their share.

\caution{A good pre-cutoff fit is evidence only if the estimator could not have fitted
any series as well. Every estimator that weights donors to match the pre-period overfits
when the donors are many for the pre-cutoff dates, and augmented SCM (\code{scm}) most of
all: its ridge features are the units' pre-cutoff outcomes, so at a pre-cutoff date the
treated unit's own value at that date is one of the features it is predicted from, and
with as many donors as dates it reproduces almost any pre-period. On the ULEZ panel of
Chapter~\ref{ch:scm} (27 pre-cutoff months), a treated series replaced by random noise is
fitted with $R^2 = -0.05$ from 5 donors, 0.69 from 20 and 0.997 from 53; classic SCM, with
40 random-walk donors over 12 months, fits a shuffled series with 0.60. So the SCM Studio
refits each design with the treated unit's pre-cutoff values shuffled in time, and reads
the pre-cutoff $R^2$ only when that shuffled series cannot also reach 0.5
(Section~\ref{sec:scmfit}). The \code{scm} estimator's donor weights balance only the
ridge residuals, so their concentration does not describe which donors the synthetic unit
is made of.}

\subsection{Inference}\label{sec:scminference}
\begin{description}
  \item[Placebo in space] Each donor in turn is treated as if it had been treated, and
        given a synthetic control built from the \emph{other} donors (the real treated
        unit is never a donor to a placebo). With the post-cutoff mean effect as the
        statistic, the permutation $p$-value is
        \begin{equation}
          p = \frac{1 + \#\{\,j : |\bar{e}_j| \ge |\bar{e}_\text{treated}|\,\}}{1 + J},
        \end{equation}
        for $J$ placebo units; the smallest it can be is $1/(J+1)$. The placebo effects
        also give a pointwise band for the treated effect.
  \item[Placebo in time] The cutoff is moved to earlier dates (every \code{placebo\_every}
        steps, keeping at least \code{min\_pre\_period} steps before it), where nothing
        happened. The true effect is ranked against the fake cutoffs' effects in the same
        way. Every fake post-period must end before the true cutoff.
  \item[Jackknife bands] The fit is repeated leaving out each donor in turn; the spread of
        the effect gives a pointwise band. A result carried by one donor shows here.
  \item[Bootstrap bands] The fit is repeated $B$ times on random subsets of the donors
        (80\,\% by default), optionally with the pre-cutoff period resampled in blocks.
  \item[Conformal interval] \code{conformal\_effect\_interval} gives a finite-sample
        interval for the average post-period effect (Chernozhukov, W\"uthrich and Zhu,
        2021); \code{rmspe\_ratio\_test} gives Abadie's post/pre RMSPE-ratio test.
\end{description}

\subsection{Staggered adoption}\label{sec:scmstaggered}
When the intervention rolls out unit by unit — the UK ULEZ 2021 expansion, staggered
driving restrictions across cities — each adopting unit has its own cutoff.
\code{run\_scm\_staggered} fits each treated unit on its own pre- and post-period, with
the \emph{other treated units excluded from every run's donor pool}, so one unit's
intervention cannot leak into another unit's synthetic control. Each unit's effect is
the post-cutoff average of its effect path,
\begin{equation}
  \bar{e}_j = \frac{1}{|T_j|} \sum_{t \ge c_j} \bigl(\text{observed}_t -
  \text{synthetic}_t\bigr),
\end{equation}
where $c_j$ is the unit's cutoff and $T_j$ its post-cutoff dates; the result is one row
per unit ($c_j$, the pre- and post-period lengths, the pre-period fit $R^2$, the effect,
and $p_j$, the unit's own placebo-in-space test at its own cutoff) plus mean and median
rows over the units' effects.

The \emph{mean} row's $p$-value is the cross-unit permutation test: each donor in turn
takes every treated unit's role — the full set of cutoffs — and its pseudo-effects are
averaged the same way the real units' are,
\begin{equation}
  p = \frac{1 + \#\{\,u : |\bar{e}_u| \ge |\bar{\bar{e}}|\,\}}{1 + J},
\end{equation}
for $J$ donors and $\bar{\bar{e}}$ the mean of the treated units' effects. The real
treated units are not permuted into the reference: their data already carry the
intervention, and in a staggered design a unit given a later cutoff than its own has its
real effect inside its new ``pre-period'', while one given an earlier cutoff has a
``post-period'' that is partly effect-free. Permuting treated identities is contaminated
by the true effect, and with few treated units the test degenerates — with two there is
one non-identity permutation, and it is contaminated. A donor-only reference is the clean
null: a donor never intervened on draws from the null distribution. Each placebo fit
needs a pool of its own, and with the default \code{scm} backend that is at least three
donors; with a smaller pool the $p$-values come back NaN while the effects are still
read.

% ============================ 4. GETTING STARTED ==============================
\chapter{Getting started}\label{ch:start}

\section{The window before a table is open}
\nm{} opens on its \tab{Data} tab with nothing loaded (Figure~\ref{fig:welcome}), the
parameter panel folded away and the Assistant (Chapter~\ref{ch:assistant}) open beside the
results. \emph{Expand Panel} on the action bar, or Ctrl+B, brings the panel out;
\tab{View $\rightarrow$ Assistant} closes the Assistant. Under
\emph{Start with a time series} are two buttons --- \emph{Open Data\ldots} and \emph{Ask
the Assistant} --- and cards for the ways in: the examples, the
five tables opened most recently (each with its folder, size and the date it last
changed; one that has since moved is greyed out), and the two places data can come from
without a file of your own, the Data Studio and the Transport Studio. A card with nothing
to list is left out. Dropping a table file (Section~\ref{sec:input}) anywhere on the
window opens it.

\screenshot{ui_welcome.png}{The window as it opens: the welcome on the \tab{Data} tab, the
action bar above the tabs (\emph{Expand Panel} first), the Assistant on the right and the
log at the bottom. The parameter panel is folded away until it is needed.}{fig:welcome}

\section{The examples}\label{sec:examples}
The examples are under \tab{File $\rightarrow$ Load Example} and on the
welcome's \emph{Examples} card; each opens with its target and predictors set as its
source sets them.
\begin{description}
  \item[London Marylebone Road \NOtwo{} (2020)] The \nm{} paper's de-weathering case
        study: hourly \NOtwo{} at a kerbside site with ERA5 meteorology, January to August
        2020, the UK's COVID-19 lockdown in it (\code{load\_my1()} below). This manual's
        main-window figures are made from it.
  \item[London Marylebone Road \PMtwo{} (2020)] The paper's transport-aware case: hourly
        \PMtwo{} at the same site, its own wind and temperature with ERA5's other fields
        (\code{load\_my1\_pm25()}).
  \item[Synthetic site, one year hourly] A synthetic year of hourly data generated on the fly: \PMtwo{}, PM\textsubscript{10},
        \NOtwo{}, O\textsubscript{3}, SO\textsubscript{2} and CO with ERA5-style
        meteorology (\code{u10}, \code{v10}, \code{t2m}, \code{d2m}, \code{blh},
        \code{sp}, \code{ssrd}, \code{tcc}, \code{tp}, \code{rh2m}, \code{ws},
        \code{wd}). The pollutants mix a seasonal cycle, a rush-hour cycle and a
        dependence on wind speed, so every step has something to show. The target is set
        to \PMtwo{} and the meteorology ticked.
  \item[Synthetic control panel] Opens the SCM Studio with its own example; the SCM
        Studio's \tab{File $\rightarrow$ Load Example} has two (Chapter~\ref{ch:scm}): the
        paper's London ULEZ panel and a synthetic weekly panel of 31 units, in which the
        unit ``2+26 cities'' has its SO\textsubscript{2} cut by about 35\,\% from 2015-10-23
        and every other unit is a clean donor.
\end{description}
The real records of the \nm{} model-description paper ship with the package; in a script
they are loaded with:
\begin{center}\small
\begin{tabular}{@{}lp{7.0cm}p{3.2cm}@{}}
\toprule
Loader & Contents & Used for \\
\midrule
\code{load\_my1()}      & Hourly \NOtwo{} + ERA5 met, London Marylebone Road, Jan--Aug 2020 & de-weathering \\
\code{load\_scm()}      & Monthly de-weathered \NOtwo{}, 104 UK sites, 2016--2021 & synthetic control (ULEZ) \\
\code{load\_my1\_pm25()} & Hourly \PMtwo{} + met at MY1, Jan--Aug 2020 & transport-aware case \\
\code{load\_traj\_my1()} & 6-hourly HYSPLIT trajectory features arriving at MY1 & transport-aware case \\
\code{example\_traj\_dir()} & Two days of raw HYSPLIT \code{tdump} files & Transport Studio \\
\bottomrule
\end{tabular}
\end{center}
Sources: UK AURN (Defra, Open Government Licence v3.0); ERA5 (Copernicus Climate Change
Service); HYSPLIT (NOAA ARL).

\section{A first analysis}
\begin{enumerate}\nosep
  \item \tab{Load data.} \emph{Open Data\ldots} and choose a table with its times, the
        pollutant and its meteorology (Chapter~\ref{ch:data}); set the \emph{Time zone} the
        site keeps (Section~\ref{sec:clock}). Bring the
        parameter panel out with \emph{Expand Panel} (Ctrl+B), or ask the Assistant.
  \item \tab{Choose the columns.} Check the target and tick the predictors
        (Section~\ref{sec:columns}).
  \item \tab{Train} and read the \tab{Model} tab: is the held-out fit good
        enough to normalise with? (Chapter~\ref{ch:train})
  \item \tab{Normalise} and compare the normalised series with the observed
        (Chapter~\ref{ch:normalise}).
  \item \tab{Decompose}, \tab{Rolling} and \tab{Partial dependence} as the question
        needs (Chapters~\ref{ch:decompose}--\ref{ch:pdp}); the \tab{Trend} test says
        whether the series the decomposition just drew is rising, falling or flat, and
        the \tab{Uncertainty} tab where the band's spread comes from (Section~\ref{sec:uncertainty}).
  \item \tab{Keep the result:} export the tables, an HTML report and the session as a
        script (Chapter~\ref{ch:export}).
\end{enumerate}

% ============================ 5. THE INTERFACE ================================
\chapter{The interface}\label{ch:ui}

The main window has a \emph{parameter panel} on the left and \emph{result tabs} on the
right. The panel has a page for each step of the workflow, chosen from the bar at its top
--- \tab{Data}, \tab{Model}, \tab{Normalise}, \tab{Decompose}, \tab{Rolling} and
\tab{PDP} --- and each step has its own \emph{Run} button beside the settings that govern
it. Results
land in the tab of the same step, which comes to the front when the step finishes. The
panel starts folded away, leaving the width to the results and the Assistant;
\emph{Collapse / Expand Panel} on the action bar (Ctrl+B) folds it and brings it back. The
window can be narrowed to about 740 pixels with the panel folded: the action bar then
wraps onto a second row.

\section{The parameter panel}
Its pages, and what each holds (a page's first box goes untitled, the tab naming it):
\begin{description}\nosep
  \item[Data] \emph{Open Data\ldots}, the file's name with its rows and columns, the
        \emph{Time zone} its times are kept on and how they were read
        (Section~\ref{sec:input}); then \emph{Columns}: the target and the predictors
        (Section~\ref{sec:columns}).
  \item[Model] Backend, split, training fraction, time budget (or search trials),
        estimators, seed; for Chronos-2 the device and fine-tuning instead
        (Chapter~\ref{ch:train}).
  \item[Normalise] The variables resampled, the number of samples, CPU cores and the
        quantile band (Chapter~\ref{ch:normalise}).
  \item[Decompose] Emission or meteorology; for meteorology the attribution and the
        groups (Chapter~\ref{ch:decompose}). Under it \emph{Multi-scale decomposition}:
        three window widths (Section~\ref{sec:multiscaleui}); \emph{Trend}: the
        series to test (a part of the decomposition, its separated long-term trend, or
        the observed target), its seasonal period, the number of segments and how
        many bootstrap resamples; \emph{Events \& anomalies}: the series to score,
        the method, the threshold and the shortest flagged stretch
        (Section~\ref{sec:eventsui}); and \emph{Uncertainty decomposition}: the number
        of ensemble members, the met draws they all share, and the members' backend
        (Section~\ref{sec:uncertainty}).
  \item[Rolling] Window, step and samples (Section~\ref{sec:rollingui}).
  \item[PDP] Variables and grid points (Chapter~\ref{ch:pdp}).
\end{description}
The panel is as wide as its widest page needs. A configuration of the panel is saved and
read from the \tab{File} menu (Section~\ref{sec:config}).

\section{Result tabs}
\begin{center}\small
\begin{tabular}{@{}lll@{}}
\toprule
Tab & Shows & Filled by \\
\midrule
Data        & overview, cycles, wind and pollution roses, statistics; rows left out & loading \\
Model       & verdict, statistics table, parity plot, feature importance, walk-forward CV & Train model; Walk-forward CV \\
Normalise   & observed and normalised series, the quantile band, their cycles & Run normalisation \\
Decompose   & the emission or meteorological components, the cycles & Run decomposition \\
Rolling     & the mean of the rolling windows, with their spread & Run rolling normalisation \\
Multi-scale & the scale-space cascade and the energy in each band & Multi-scale \\
Trend       & the tested series with its Theil--Sen line, or its segments and marked break, the verdict and the p value & Trend analysis \\
Events      & the anomaly-score strip with the flagged windows shaded, and the event table & Find events \\
PDP         & one partial-dependence panel per variable & Compute PDP \\
Uncertainty & the per-date spread stacked by source (weather, model, interaction) and the median's shares & Uncertainty decomposition \\
\bottomrule
\end{tabular}
\end{center}
A tab with a judgement to make opens with a \emph{verdict banner} in green (good), amber
(caution) or red (not usable), followed by plain lines of what was run and, where there
is one, a compact table. Every figure has a toolbar of small icons in the window's
colours --- home, back, forward, pan, zoom, subplot spacing, axes and curve options,
save --- then \emph{For print\ldots} (Section~\ref{sec:print}) and, at its right end, the
value under the pointer. A figure is saved on a white page with dark ink whatever the
window's appearance. Every figure is drawn as journals ask for one: 7 to 8.5~pt text, thin
axes with ticks outward, legends without a frame, colours a colour-blind reader can tell
apart, formulas with their subscripts (\NOtwo{}, \PMtwo{}), and a date axis written once
at each level --- the year once, then the months --- with panels stacked on one time axis
dated on the bottom panel only.

\paragraph{Resolution.} The time series on the \tab{Data}, \tab{Normalise},
\tab{Decompose}, \tab{Rolling} and \tab{Multi-scale} tabs have a \emph{Resolution} selector above the
figure: \emph{As recorded}, \emph{Hourly}, \emph{Daily}, \emph{Weekly}, \emph{Monthly}
or \emph{Yearly}. A coarser step redraws the figure from the means of each period, which
makes a long hourly record readable; steps finer than the record's own are greyed out.
A record finer than daily and longer than 2,000 points opens as daily means --- months of
hours drawn as they are are a solid band --- and the choice made on a tab is kept for its
next result. Only the display changes: the result, its table and its export stay at the
record's resolution. A series is drawn broken where the record has a gap (an interval
over five times its usual step), not joined across it by a straight line. The $\sigma$ a
figure quotes are the record's, whatever step it is drawn at.

When both an emission and a meteorological decomposition have been run on the model, the
\tab{Decompose} tab keeps both: a \emph{Method} selector beside \emph{Resolution} switches
between them, and a saved session brings both back.

\section{The action bar}
Above the tabs: \emph{Collapse / Expand Panel}; the studios side by side ---
\emph{Get data} (the Data Studio, Section~\ref{sec:datastudio}), \emph{Transport
Studio} (Section~\ref{sec:transport}) and \emph{SCM Studio} (Chapter~\ref{ch:scm}); then
what is done with a result --- \emph{Export result CSV} (the current tab's table) and
\emph{HTML report} (the current tab's result as a report with its provenance).

\section{Run buttons}
\begin{center}\small
\begin{tabular}{@{}lp{5.0cm}l@{}}
\toprule
Button & Needs & Runs \\
\midrule
Train model / Load Chronos-2 & a table, one predictor ticked & \code{build\_model} \\
Fine-tune Chronos-2   & Chronos-2 loaded         & \code{finetune} \\
Run normalisation     & a model, one variable to resample & \code{normalise} \\
Run decomposition     & a model                  & \code{decompose} \\
Run rolling normalisation & an AutoML model      & \code{rolling} \\
Multi-scale decomposition & an AutoML model, a normalisation & \code{rolling} $\times 3$ \\
Trend analysis      & an AutoML model (a decomposition, where it tests one) & \code{trend} (1 segment), \code{detect\_trend\_breaks} (more) \\
Find events         & a model (a decomposition, where it scores a part) & \code{anomaly\_scores}, \code{detect\_events} \\
Compute PDP           & an AutoML model, one variable ticked & \code{pdp} \\
Uncertainty decomposition & an AutoML model & \code{decompose\_uncertainty} \\
\bottomrule
\end{tabular}
\end{center}
A button is enabled only when it can act, and its tooltip says what is missing when it
cannot. One task runs at a time, on a background thread, so the window stays responsive;
while it runs the run buttons are disabled, \emph{Cancel} is enabled, and the status bar
names the task with the time it has taken so far. Where the task says how far it has come
--- the rolling windows, the decomposition's steps, LightGBM's trials, the receptor times
HYSPLIT runs, FLAML's time budget --- the progress bar fills with it, and the line gives
the share done and, once there is enough to go on, the time left; elsewhere the bar is a
busy bar. When the task ends, the line its step writes says how long it took. Cancel \emph{abandons} the task: the window is
free at once, but the computation cannot be interrupted, so it finishes in the background
and its result is discarded. Each step's call is added to the command log when it
succeeds (Section~\ref{sec:commandlog}).

A step asked for before what it needs --- normalising before a model is trained, a
multi-scale run before the normalisation --- is answered by a \emph{notice} above the results, not a
dialog: it says what is missing and carries the step that answers it on a button
(\emph{Train model}, \emph{Run normalisation}); \(\times\) puts it away, and it goes by itself
once the step has been taken. At the right of the status bar a light per step --- Data,
Model, Normalised, Decomposed, Multi-scale, Trend, Events, Uncertainty, Rolling, PDP ---
turns green when that step has a result.

\section{Menus and shortcuts}
\begin{center}\small
\begin{tabular}{@{}lll@{}}
\toprule
Menu & Command & Shortcut \\
\midrule
File     & Open Data\ldots                      & Ctrl+O \\
         & Load Example (the examples, SCM's)   & \\
         & Get Data (AQ networks + met)\ldots & Ctrl+D \\
         & Open Recent (eight tables)           & \\
         & Save Results\ldots{} / Load Results\ldots & Ctrl+S \\
         & Reopen Last Session                  & \\
         & Load Session Script\ldots            & \\
         & Methods Paragraph\ldots              & \\
         & Export Current Result\ldots          & Ctrl+E \\
         & Export HTML Report: English, Chinese & \\
         & Save Session as Script\ldots         & \\
         & Save Config\ldots{} / Load Config\ldots & \\
         & Run History\ldots                    & \\
View     & Command Palette\ldots                & Ctrl+K \\
         & Collapse / Expand Panel              & Ctrl+B \\
         & Log\ldots, Command Log\ldots       & \\
         & Assistant                            & Ctrl+Shift+A \\
         & Appearance: Follow System, Light, Dark & \\
         & Language: Follow System, English, Chinese & \\
Analysis & Synthetic Control (SCM) Studio\ldots & Ctrl+Shift+S \\
         & Polar Plot\ldots                     & \\
Help     & User Manual                          & F1 \\
         & Keyboard Shortcuts\ldots             & \\
         & Documentation, About Normet          & \\
\bottomrule
\end{tabular}
\end{center}
On macOS, Ctrl is the Command key. The \emph{command palette} (Ctrl+K) lists every
menu command and every result tab; type a few letters of one and press Return.
\tab{Help $\rightarrow$ Keyboard Shortcuts} lists every shortcut, read from the menus;
\tab{Help $\rightarrow$ User Manual} (F1) opens this manual in the window's language.

\section{The log and the Assistant}
\tab{View $\rightarrow$ Log\ldots} opens \nm{}'s log in a window of its own, which can
stay open beside the main window: what each step did, how many rows, FLAML's chosen
estimator and its hyperparameters, how many windows a rolling run made.
\tab{View $\rightarrow$ Command Log\ldots} opens the calls that repeat the session
(Section~\ref{sec:commandlog}). When a task fails, its full traceback is written to the
Log and its last line shown in a dialog. The \emph{Assistant} takes the right of the
window, its whole height (Chapter~\ref{ch:assistant}).

\section{Appearance and language}\label{sec:appearance}
The window and its figures follow the system's light or dark appearance;
\tab{View $\rightarrow$ Appearance} pins one. A figure on a dark window is recoloured
for display only --- its faces, text, ticks and grid; colours that carry data are left
alone --- and saved in its light colours.

The window speaks English or Chinese: the system's language by default, or the one
\tab{View $\rightarrow$ Language} chooses, kept for the next session. The switch is
immediate --- menus, buttons, tabs, tooltips, verdicts and the figures' titles, labels and
month names are written again in the language chosen, and a figure is saved as it is
shown. The data's own words (column names, site codes), the log, the command log's code
and the Assistant's conversation stay as they were written. What the system draws
itself stays in the system's language whichever the window is written in: the
open and save panels, their titles too, and, on a Mac, the items macOS adds to the menus
(\emph{AutoFill}, \emph{Start Dictation}, \emph{Emoji \& Symbols}). An HTML
report (Section~\ref{sec:commandlog}) is written in the language picked as it is exported
--- \tab{File $\rightarrow$ Export HTML Report}, or the \emph{HTML report} button, lists
each language named in itself --- its headings, labels and figure; the language is kept
for the next report, and is the system's until one is picked.

% ============================ 6. COLUMNS AND TRAINING =========================
\chapter{Training a model}\label{ch:train}

\section{The table and its columns}\label{sec:columns}
A table opened in the main window needs a column called \code{date} (any case) whose
values parse as dates; every other column numeric is offered as a target and as a
predictor. The target is preset to the first column named like a pollutant (\PMtwo{},
PM10, \NOtwo{}, O3, SO2, CO, NOx), and the predictors to the columns recognised as
meteorology:
\begin{center}\small
\code{u10 v10 d2m t2m blh sp ssrd tcc tp rh2m ws wd temp rh pressure}
\end{center}
(matched without regard to case). \emph{All} ticks every predictor --- and, once all are
ticked, unticks every one --- and \emph{Met only} ticks the meteorology alone. The target is never a predictor, even when ticked. The four calendar
variables are added to every model automatically and are not listed.

\note{A meteorological column under another name is not recognised and must be ticked by
hand --- in MY1, \code{atmos\_pres}. The same applies to the \tab{Normalise} page's \emph{Met only}
(Chapter~\ref{ch:normalise}).}

\screenshot{ui_data.png}{MY1 opened from \tab{File $\rightarrow$ Load Example}: 5,793
hourly rows of 11 columns from 2020-01-01 to 2020-08-31, with no missing values. The
example sets the paper's target, \NOtwo{}, and its nine meteorological predictors
(\code{atmos\_pres} among them, a name the window does not recognise). The \tab{Overview}
page draws every column as a heatmap, each scaled to its own range, above the series
picked (here \NOtwo{}, daily).}{fig:data}

The \tab{Data} tab has four pages. \tab{Overview} shows the record at a glance: drag across
the plot to pick a span, or click it to pick the point under the cursor (the period it
stands for, where the series is averaged), and the rows are listed below the plot; with
Shift held (or Ctrl, or Cmd) a drag or a click adds to the pick, a click on a picked point
drops it, and Escape clears them. One line under the plot says the pick --- its rows,
and how many of them are left out --- at its left, and the time, the variable and the
value under the cursor at its right. In the list, rows are selected
as in a spreadsheet --- Ctrl- or Cmd-click for rows apart, Shift-click or a drag down the
rows for a run --- and right-clicked, on them or on their numbers, to exclude or restore
them (a row right-clicked that is not selected is selected alone first). A right-click on
the plot offers the same for every row picked, a point outside the pick picked first;
\emph{Exclude picked}
and \emph{Restore picked}, on the plot's \emph{Resolution} row, do the same for the rows
selected, or every row picked when none is. The list stays, its excluded rows in grey, to
go on from. The \emph{Exclude} box beside them takes spans typed one to a line
(\code{2020-03-23 .. 2020-04-05}, or a single date) and \emph{Apply}. Rows left out
are shaded on the plot and are left out of every step after: training, normalisation,
decomposition and the rest. The Command Log writes them as \code{left\_out} and
\code{df\_kept}, and the session keeps them. \tab{Cycles} draws a column's diurnal, weekly and
seasonal cycles; \tab{Wind}, when the table has wind speed and direction, the wind rose,
a pollutant's pollution rose and its polar plot; \tab{Statistics} each column's values,
missing share, mean, standard deviation, minimum, quartiles and maximum. All four read
the rows kept.

\section{Settings}
\begin{description}\nosep
  \item[Backend] \code{flaml}, \code{lightgbm} or \code{chronos-2}
        (Section~\ref{sec:chronosui}). Chronos-2 is greyed out when its extra is not
        installed, its tooltip saying why.
  \item[Split method, Training fraction] Section~\ref{sec:split}; 0.75 by default.
  \item[Time budget] FLAML's search time, 60~s by default. For LightGBM the row becomes
        \emph{Search trials}, 20 by default; each backend's value is kept when you switch.
  \item[Estimators] (FLAML only) LightGBM alone by default; add more to widen the search.
  \item[Seed] 7654321 by default: the split and the search are seeded with it.
\end{description}

\section{Reading the Model tab}
\emph{Train model} runs \code{build\_model} --- preparation, the split and the search ---
and fills the \tab{Model} tab (Figure~\ref{fig:model}). The verdict is read from the
held-out $R^2$: 0.7 or above is a good fit, 0.5 to 0.7 moderate, below 0.5 poor. The table
is \code{modStats} for the training rows, the testing rows and all rows
(Section~\ref{sec:modstats}); the parity plot shows the held-out rows against the
prediction, square, with the 1:1 line and, in its corner, the number of rows, the slope of
predicted on observed (below 1: the extremes are pulled in) and the mean bias --- as a
density of hexagons where the rows run to thousands; the bars are the model's feature
importance (split counts for LightGBM). Under the table, one line reads the point
forecast's \emph{CRPS reference} --- $\sqrt{\pi/2}\cdot\mathrm{MAE}$ over the held-out
rows, the point forecast's score on the scale of the continuous ranked probability score.
A forecast that carries a spread scores its own CRPS, and only a score below this
reference says the spread is worth carrying.

\screenshot{ui_model.png}{MY1 trained with FLAML (LightGBM, 30~s budget, random split):
on the 1,448 held-out hours $R^2 = 0.883$ and RMSE $= 8.66$~\ugm{}, against $0.976$ on
the 4,345 training hours. Wind direction and speed, the boundary-layer height and the
trend carry the most splits.}{fig:model}

A fit on the training rows much better than on the held-out rows is expected of boosted
trees and is not by itself a problem; a poor held-out fit is. Before normalising with a
moderate or poor model, try more predictors, a longer budget or more estimators. The
normalised series inherits the model's errors.

A single random split flatters the model: every test hour's neighbours sit in the
training rows, so an autocorrelated series can be forecast locally without being
forecast well. \emph{Walk-forward CV} checks that number: the same backend retrains on
an expanding window and scores the block of hours immediately after it, five times by
default (\emph{CV folds}); it takes the training budget once per fold. The tab gains a
\emph{View} selector between the split's figures and the folds --- a bar of $R^2$ per
fold with the split's value drawn across the bars, and a table of each fold's test
window, its size and its $R^2$, RMSE and mean bias. When the split's line sits well
above the bars the verdict says so: trust the folds. For a series whose regime drifts,
where a random split samples the future as often as the past, the folds are the honest
number.

\section{Chronos-2}\label{sec:chronosui}
With \code{chronos-2} chosen the training controls are hidden --- there is nothing to
train, no budget, no split --- and the button reads \emph{Load Chronos-2}. The first load
downloads about 500~MB of weights. \emph{Device} chooses where it runs (\emph{auto}:
CUDA, else Apple Silicon, else CPU). At least one meteorological predictor must be
ticked: without covariates the model makes a plain forecast, not a de-weathering.

The \tab{Model} tab then reports the checkpoint, device, context and horizon, and the
covariate-sensitivity check (Section~\ref{sec:chronos}): the mean and largest shift of the
forecast's median when the future weather is shuffled, and that shift as a percentage of
the prediction. The verdict is green at 5\,\% or more, amber from 1 to 5\,\% (the series is
mostly carried by the target's own history) and red below 1\,\% (the model is
autoregressing, and a de-weathered series from it would be meaningless).

The table also carries a quantile backtest: the checkpoint forecasts the last horizon once
at the 10/50/90 percentile levels, the band is scored against the observed tail by CRPS
(the CDF interpolated through the three levels), and its median is scored as a point
forecast ($\sqrt{\pi/2}\cdot\mathrm{MAE}$ again). The line under the table says whether the
band beats its own median's reference --- whether the spread is worth carrying --- and the
pinball losses break the score down by level. The line names the window it speaks for: one
horizon is a quick check, not the backtest one would quote.

\emph{Fine-tune Chronos-2} adapts the loaded checkpoint to the site's own record:
\emph{lora} trains small adapters and leaves the base weights alone (it needs
\code{peft}); \emph{full} updates all 119~million parameters and overfits a single site
more easily. \emph{Steps} sets the optimiser steps (300). The tuned weights replace the
model for the session; normalise again to use them, or load Chronos-2 again to go back.
On the Chronos-2 backend, rolling, partial dependence and the emission decomposition
are disabled, and the multi-scale decomposition, built on rolling, does not apply
(Section~\ref{sec:chronos}).

% ============================ 7. NORMALISATION ================================
\chapter{Normalising}\label{ch:normalise}

\begin{description}\nosep
  \item[Resample vars] Every variable the model was trained on, the four calendar
        variables included. Ticked variables are resampled (their influence averaged
        out); unticked ones keep their observed values. After training the list is set
        to \emph{Met only} --- the usual de-weathering --- and \emph{All} ticks
        everything, which removes the trend and the seasonal and weekly cycles as well.
  \item[Samples] Monte-Carlo draws, 300 by default (8 on Chronos-2;
        Section~\ref{sec:normalise} on how many are enough).
  \item[CPU cores] \emph{auto} uses all but one.
  \item[5--95\,\% quantile band] Also return the 5th, 25th, 75th and 95th percentiles of
        the draws and shade the outer band.
\end{description}

Beneath the series, the observed and normalised series are averaged by the hour of the
day, the day of the week and the month, each mean with its 95\,\% confidence interval ---
the record's cycles, as openair's \code{timeVariation} draws them, and how much of each
the weather moved. A cycle the record cannot show (the hours of a daily record) is left
out; the cycles are averaged from the record as it is, whatever resolution the series
above is shown at. In a script: \code{nm.cycle\_profiles(result, ["observed",
"normalised"])}.

\screenshot{ui_normalise.png}{MY1 normalised with 300 draws of all nine meteorological
variables, shown as daily means. The normalised series (red) falls in the third week
of March, the lockdown, which the weather does not explain. Beneath, its cycles: the
normalised series keeps the rush hours and the weekday--weekend difference --- the
calendar is held at its observed values --- with less of the observed series' spread.}{fig:normalise}

The line above the figure gives the number of timestamps, draws and resampled variables,
and the mean of normalised $-$ observed over the record: here $-0.26$~\ugm{}, since over
eight months the weather averages out. The monthly means tell the story the figure shows:
\begin{center}\small
\begin{tabular}{@{}lrrrrrrrr@{}}
\toprule
\NOtwo{} (\ugm{}) & Jan & Feb & Mar & Apr & May & Jun & Jul & Aug \\
\midrule
observed   & 60.3 & 65.2 & 51.7 & 34.5 & 30.7 & 35.1 & 31.6 & 36.6 \\
normalised & 59.7 & 64.7 & 54.0 & 34.7 & 32.1 & 32.5 & 32.1 & 33.8 \\
\bottomrule
\end{tabular}
\end{center}
From February to April the normalised series falls by 30.0~\ugm{} (46\,\%), against
30.7~\ugm{} observed: almost none of the spring drop was the weather. In March the weather
was cleaner than usual (observed 2.3 below normalised), in June and August dirtier.

\note{On Chronos-2 the first \code{context\_length} rows are seeded with the observed
values; they are shaded grey on the figure, the line under it says how many there are,
and the mean is taken over the rest.}

\caution{The normalised series is the model's expectation under typical weather. A model
that has not learnt the weather's effect --- a poor fit, or meteorology missing from the
predictors --- returns something close to the observed series, which looks like a
de-weathered series with nothing to remove. Read the \tab{Model} tab first.}

% ============================ 8. DECOMPOSITION ================================
\chapter{Decomposing}\label{ch:decompose}

\section{By time scale: the emission decomposition}
With \emph{Method} set to \emph{emission}, the decomposition runs the nested normalisation of
Section~\ref{sec:emission} and plots the normalised series with its trend, seasonal,
weekly, diurnal and short-term components, each with its standard deviation $\sigma$
(Figure~\ref{fig:emission}). The line above the figure names the component with the
largest $\sigma$. On a record finer than daily, the diurnal and weekly components are
drawn on the right as their cycles --- the component's mean by hour of day and by day of
week, shaded between its 25th and 75th percentiles --- rather than as months of
oscillation.

\screenshot{ui_decompose_emission.png}{The emission decomposition of MY1 (100 draws per
normalisation; baseline 43.1~\ugm{}). The lockdown appears as a step in both the trend
($\sigma = 9.56$) and the seasonal component ($\sigma = 4.74$): a one-off event is not a
season, and the record is too short to tell them apart. The diurnal cycle (right) is the
largest ($\sigma = 9.92$): low before 05:00, high through the working day, its spread wide
because it shrank after the lockdown, as rush-hour traffic did; the weekly cycle
($\sigma = 3.36$) dips at the weekend.}{fig:emission}

\section{By variable: the meteorological decomposition}
With \emph{Method} set to \emph{meteorology}, the decomposition splits observed $-$ normalised among
the meteorological variables (Section~\ref{sec:metdecomp}). \emph{Attribution} is enabled and \emph{Groups} appears
(it is hidden for \emph{emission}):
\begin{description}\nosep
  \item[Attribution] \emph{sequential}, \emph{shapley}, or \emph{auto} (Shapley when groups
        are given, else sequential).
  \item[Groups] A table of the meteorological variables the model was trained on, each
        with its group: \emph{Wind}, \emph{Temperature and humidity}, \emph{Boundary
        layer and pressure}, \emph{Sky and rain}, or a name typed in. \emph{Suggest
        groups} puts each recognised variable in its group; \emph{Each alone} clears
        them. A variable left without a group is a contribution of its own. A group's
        rows are tinted in one colour, and the line under the table says how many
        contributions the grouping makes.
\end{description}
The figure shows met\_total and one panel per variable or group, with a ranking of their
$\sigma$ beside \code{met\_noise}, the part the model does not reproduce
(Figure~\ref{fig:metdecomp}).

\screenshot{ui_decompose_met.png}{MY1's weather influence (observed $-$ normalised,
$\sigma = 16.79$~\ugm{}) split among four groups by exact Shapley values (16 coalitions):
wind (\code{ws}, \code{wd}; $\sigma = 9.92$), thermodynamics (\code{temp}, \code{RH};
5.24), mixing (\code{blh}, \code{atmos\_pres}; 4.89) and sky (\code{tcc}, \code{tp},
\code{ssrd}; 3.44).}{fig:metdecomp}

The decomposition runs $k+1$ normalisations sequentially or $2^k$ by Shapley, each of
\emph{Samples} draws, so its cost grows quickly with the number of variables. Ten
variables by Shapley is 1024 normalisations: group them. Unlike the normalisation, the
meteorological decomposition always resamples every meteorological variable trained on;
it does not follow the \tab{Normalise} page's list.

\section{Trend}\label{sec:trend}
The question a decomposition invites: is the series it just drew rising, falling or flat,
and is the rise real? The \emph{Trend analysis} box under \emph{Decompose} runs
Mann--Kendall and Theil--Sen on the chosen series and, the record carrying the
autocorrelation the correction exists for (MY1's emission part: AR(1) residual acf 0.94),
replaces the slope's classical error bar with a moving-block bootstrap. It runs again
each time an emission decomposition ends where \emph{Run after each emission
decomposition} is ticked (unticked by default).

\emph{Series} names which one is tested: the decomposition's emission part
(the default: what the weather does not explain), its \emph{separated long-term
trend} --- the long-term component Section~\ref{sec:emission} drew, cycles already
separated into their own components, so the cleaner input to a trend analysis --- its
weather part, or the observed target itself. \emph{Segments} names how many:
one (the default) is the single-slope test just described; two to four run the
\emph{segmented fit} instead --- \code{detect\_trend\_breaks} on the chosen
series. It splits the record where the trend breaks, if it does, each segment
reported with its own Theil--Sen slope and interval, and the break accepted only
under the two-part gate (Section~\ref{sec:emission} for the components it tests on):
the BIC must improve, and the SSE improvement of the final fit over the single
slope must beat the same search re-run on block-bootstrapped resamples (the p to
quote). A record that fails a gate is reported as one regime, with the best
candidate cut and its p kept so the tab can say what was looked at. Raising
\emph{Segments} above one moves a still-default \emph{emission part} choice to
the long-term trend (the clean input); the choice stays until changed.

The \tab{Trend} tab (Figure~\ref{fig:trend}) for one segment shows the series
with its Theil--Sen line and a verdict that says which of the two p values to
quote: the classical Mann--Kendall one is the anti-conservative one on an
autocorrelated record, and the record is autocorrelated. For two or more
segments it shows the series with its \emph{piecewise} Theil--Sen lines and a
marked break, a verdict that says whether the search found a break (where, and
the block-bootstrap p to quote) or that the gates refused one, per-segment rows
in the table (start, end, n, slope, 95\,\% interval, per year), and the same
CSV export and command-log entry as the one-segment run
(\code{trend\_breaks\_call} in the script, not \code{trend\_call})
(Figure~\ref{fig:trend_breaks}).

\screenshot{ui_trend.png}{The trend analysis of MY1's decomposition's emission part
(Section~\ref{sec:emission}): the series with its Theil--Sen line and the two p values ---
the classical Mann--Kendall one and, the record being autocorrelated, the
block-bootstrap one --- with the verdict above the figure saying which to quote.}{fig:trend}

\screenshot{ui_trend_breaks.png}{The segmented fit (Segments 2) on MY1's separated
long-term trend: the search found the break at 2020-03-18 19:00 --- the 2020-03-23
lockdown's lead-in --- the single-slope SSE down 88.7\,\% with the block-bootstrap
p of the search 0.001, the one to quote. Segment 1 runs +0.0812 per day
(95\,\% [0.0761, 0.0863]); segment 2 flattens to $-$0.0121 per day
(95\,\% [$-$0.0134, $-$0.0108]). The piecewise Theil--Sen lines and the marked
break are in the figure; the per-segment table and the verdict above it are the
Trend tab's.}{fig:trend_breaks}

\section{Events \& anomalies}\label{sec:eventsui}
The question beside the trend: not where the record is heading but what happened
in it. The \emph{Events \& anomalies} box under \emph{Trend} scores every
timestamp for how far it strays from the record's normal behaviour and flags the
consecutive stretches above a threshold --- the smoke episodes, the release
events, the days the instruments misbehaved (\code{anomaly\_scores},
\code{detect\_events}).

\emph{Series} is the trend box's choice: a decomposition's emission part
(events the weather did not cause), its separated long-term trend, its weather
part, or the observed target. \emph{Method} chooses the score: \code{iqr},
$|x-\text{median}|/\mathrm{IQR}$, the robust univariate score with
\emph{Threshold k} in multiples of it; \code{isolation}, an IsolationForest over
the value and its lags, which catches shapes a univariate score cannot --- a
spike train, a level shift --- and scores against the forest's own boundary, so
the threshold row goes grey; \code{stl\_residual}, the score of the
seasonal-decomposition residual, what stands out against the record's own cycle
rather than against its whole distribution. \emph{Min length} drops stretches
shorter than it, in steps of the record: 1 keeps single-hour spikes, 3 reads as
episodes rather than instrument flickers.

The \tab{Events} tab draws the score strip --- every timestamp's score with each
flagged window shaded, the threshold as a line where one applies, so the
shading is simply where the curve crosses it --- and the table of windows:
start, end, steps, maximum and mean score, ranked by the maximum. On an
autocorrelated hourly record even $k=3$ flags the tails of ordinary weather:
raise the threshold, or the minimum length, until only the episodes remain.
Both calls are deterministic under the seed, so the table's rows and the
strip's shading are the same stretches, hour for hour.

\screenshot{ui_events.png}{The \tab{Events} tab on MY1's emission part
(Section~\ref{sec:emission}): the iqr score of every timestamp, the $k=2$
threshold as a line --- the box's default $k=3$ flags no window on this
de-weathered series, so 2 is the first value that does --- and the 26 flagged
windows shaded, the table of them ranked by the maximum score above the figure,
the largest scoring 2.36. The shading is simply where the curve crosses the
line; at a low threshold the ordinary-weather tails are what get flagged, and
the top of the table is the 2020-03-10 07:00--18:00 stretch, a full day.}{fig:events}

\section{Uncertainty decomposition}\label{sec:uncertainty}
The question a de-weathered estimate invites: of the band around an effect, how
much is the met resampling and how much is the model form? The \emph{Uncertainty
decomposition} box under \emph{Events \& anomalies} answers by attribution, not by another
Monte-Carlo run. It trains \emph{Models} (default 4) members, each on its own
seed and on the table as loaded, from the trained variable set; every member then
predicts the same \emph{Samples} (default 300, the variables ticked on the
\tab{Normalise} page) met resamples, the draws seeded once by the window's seed
so that draw $j$ is identical across members. Each timestamp's
(members $\times$ draws) prediction matrix splits its variance orthogonally into
the between-draw term --- the weather variance --- the between-member term, the
model-form variance, and the interaction: what a member reacts to in a draw that
the ensemble on average does not. The three summands sum to the total per-date
variance exactly --- an identity the package pins in a test --- so the
attribution cannot quietly stop adding up. The \emph{Member backend} chooses how each member
is trained: \code{lightgbm} (the default: a fast random search with the Model
page's budget as its trials) or \code{flaml} (an AutoML search per member, at
that member's budget each).

The \tab{Uncertainty} tab (Figure~\ref{fig:uncertainty}) stacks the per-date
variance, weather under model under interaction, the line over the stack the
total, and says the headline above the figure: the median share of the total by
source. For MY1, three members and 50 shared draws under the figure's settings,
the median spread is 96\,\% met resampling, 3\,\% the interaction and 0\,\%
model form --- the fitted model is nearly certain about what the weather does,
and the band an effect is tested against is the weather's, not the model's. The
table gives each source's median variance and median share; the shares are
medians across the record, so a single extreme day does not re-weight them. The
run is added to the command log as the \code{decompose\_uncertainty} call the
window made, and a saved session brings the run back whole, its figure re-drawn.
The cost is \emph{Models} fits and \emph{Models} normalisations of
\emph{Samples} draws --- the same number of predictions as one
\code{do\_all\_unc} run, and no more; a run of MY1's settings takes a minute or
two, not an hour.

\screenshot{ui_uncertainty.png}{The attribution of MY1's per-date spread: the
stacked bands (met resampling, model form, the interaction) total 144.3
(\ugm{})$^{2}$ of variance in the median, and the verdict above says where it
goes --- 96\,\% met resampling, 3\,\% interaction, 0\,\% model form (three
members sharing 50 met draws, lightgbm members). An effect smaller than that band
is not separated from the weather by the model form at all.}{fig:uncertainty}

% ============================ 9. ROLLING ======================================
\chapter{Rolling and multi-scale}\label{ch:rolling}

\section{Rolling normalisation}\label{sec:rollingui}
\emph{Window} (14~d), \emph{Step} (7~d) and \emph{Samples} (100) set the windows of
Section~\ref{sec:rolling}; the variables resampled are the \tab{Normalise} page's. The tab shows the mean of
the overlapping windows at each time, with $\pm 1$ standard deviation across them where
two or more overlap (Figure~\ref{fig:rolling}). The result table, exported, has one column
per window.

\screenshot{ui_rolling.png}{Rolling normalisation of MY1: 33 fortnightly windows started
weekly, 50 draws each. Each window removes only the weather that varies within its own
fortnight.}{fig:rolling}

\section{Multi-scale decomposition}\label{sec:multiscaleui}
The multi-scale box runs rolling normalisations at three widths --- \emph{Fast} (14~d),
\emph{Meso} (90~d) and \emph{Slow} (365~d) by default --- and differences them against each
other and against the normalised series (Section~\ref{sec:multiscale}); so the
normalisation must have been run first. It uses the \tab{Rolling} page's \emph{Samples} and costs about three rolling runs. The
verdict names the band with the most energy ($\sigma$); a band that could not be computed
is named in the notes, with why.

\screenshot{ui_multiscale.png}{Multi-scale decomposition of MY1 at 14, 60 and 120~days
(the 365-day default is longer than this eight-month record): 77, 13 and 5 windows. The
synoptic band carries the most ($D_\text{fast}$, $\sigma \approx 3.0$), then the intra-seasonal
and the slow.}{fig:multiscale}

% ============================ 10. PDP =========================================
\chapter{Partial dependence}\label{ch:pdp}

After training, the variable list holds every predictor and the calendar variables, with
the first four ticked. \emph{Time variables} and \emph{Met only} toggle those groups;
\emph{Grid points} (50) sets the resolution. Each panel shows the curve
(\code{pdp\_mean}) with a band of $\pm$ \code{pdp\_std}, the spread of the predictions
across the training rows at that value (Section~\ref{sec:pdptheory}). The short ticks
along the bottom are the variable's deciles in the training data: where they are far
apart --- a curve's ends --- the data are few. A variable of a few whole values (the hour,
the weekday) is drawn as steps, a point on each value.

\screenshot{ui_pdp.png}{Partial dependence of the MY1 model: \NOtwo{} falls as wind
speed rises and as the boundary layer deepens from its shallowest, rises above about
20\,$^\circ$C, and follows the traffic day --- low before 05:00, high from 07:00 to
17:00.}{fig:pdp}

The bands are wide because the other predictors still vary across the rows; they are not
a confidence interval for the curve.

% ============================ 11. SCM STUDIO ==================================
\chapter{Synthetic control (SCM Studio)}\label{ch:scm}

The synthetic-control workflow takes a different input --- a long panel of many units
--- so it has a window of its own, opened from \emph{SCM Studio} on the action bar,
\tab{Analysis $\rightarrow$ Synthetic Control (SCM) Studio} (Ctrl+Shift+S) or the
welcome's example card. It is laid out as the main window is: the panel on the left,
its pages chosen from a bar --- \tab{Data} (the panel data and its columns),
\tab{Design}, \tab{Estimator}, \tab{Inference} --- result tabs on the right
(\tab{Panel}, \tab{Fit}, \tab{Fit Checks}, \tab{Placebo Space}, \tab{Placebo Time},
\tab{Uncertainty}, \tab{All Units}, \tab{Similarity}), and \emph{Export current tab
data} above them. The
method is described in Section~\ref{sec:scmtheory}.

This chapter's figures reproduce the ULEZ case study of the \nm{} paper: the London Ultra
Low Emission Zone, launched on 8~April 2019, evaluated at Marylebone Road (MY1) on
monthly de-weathered \NOtwo{}, with traffic sites elsewhere in the UK as donors. To keep
COVID-19 out of the post-period the panel ends in February~2020, and only traffic sites
with a complete record are kept: 53 donors over 38 months. The SCM Studio's
\tab{File $\rightarrow$ Load Example $\rightarrow$ London ULEZ} sets this design up; in a
script:
\begin{lstlisting}[language=Python]
import normet as nm
scm_df = nm.datasets.load_scm()
df = scm_df.query("'2017-01-01' <= date < '2020-03-01'")
df = df[df["type"].str.contains("Traffic")]
counts = df.dropna(subset=["NO2_dw"]).groupby("code").size()
donors = [c for c in counts[counts == df["date"].nunique()].index if c != "MY1"]
df = df[df["code"].isin(donors + ["MY1"])]
df.to_csv("ulez_panel.csv", index=False)
\end{lstlisting}

\section{Panel data and columns}
\emph{Open Data\ldots} (also \tab{File $\rightarrow$ Open Data\ldots}) reads a long
table, in any of the formats of Section~\ref{sec:input}: one row per unit and date. \emph{Use
main-window data} takes the main window's table, when it has a unit column. The
\emph{Date} column is guessed from its name, the \emph{Unit} column as the first
non-numeric column with between 3 and a quarter as many distinct values as rows, and the
\emph{Outcome} as a numeric column whose name ends in \code{wn} (the de-weathered series
of the paper's tables), else the first numeric one. Check all three.

\section{Design}
\begin{description}\nosep
  \item[Treated unit] The unit the intervention reached.
  \item[Cutoff date] The first date of the post-period. The month containing the
        intervention is a judgement: here 2019-04-01, the ULEZ's launch month.
  \item[Donor pool] Every other unit, ticked. Untick any unit the intervention could
        have reached, directly or by spill-over --- a donor that was itself treated
        biases the counterfactual towards the treated unit.
  \item[Rank donors by similarity] (needs the \code{foundation} extra) Embeds every
        unit's pre-cutoff record with Chronos-2 and ranks the others by the cosine
        similarity of their dynamics to the treated unit's, on the \tab{Similarity} tab;
        \emph{Tick the top N} and \emph{Apply} make those the pool.
\end{description}
The \tab{Panel} tab plots the outcome of every unit with the treated one in red, and
checks the panel: amber when more than 10\,\% of the date $\times$ unit cells are missing
or there are fewer than 10 pre-cutoff dates, green otherwise.

\screenshot{ui_scm_panel.png}{The ULEZ panel: MY1 (red) against the 53 traffic-site donors,
monthly de-weathered \NOtwo{}, January~2017 to February~2020. MY1 is the highest of
them throughout.}{fig:scmpanel}

\section{Fitting}\label{sec:scmfit}
\emph{Backend} chooses the estimator (Section~\ref{sec:scmtheory}); \emph{ML backend} and
\emph{Time budget} (15~s) apply to \code{mlscm}, \emph{Draws} (1000) to
\code{bayesian}. \emph{Run synthetic control} fills the \tab{Fit} tab: observed and
synthetic, the effect, the largest donor weights, and the diagnostics table.

For \code{scm}, \code{abadie}, \code{did}, \code{mcnnm} and \code{robust} the fit is
followed by one more, of the same design with the treated unit's pre-cutoff values
shuffled in time: a series with the treated unit's level and spread but no relation to the
donors (Section~\ref{sec:scmdiag}). Its pre-cutoff $R^2$ is the table's
\code{pre\_r2, shuffled} row, and it decides how the verdict reads:
\begin{itemize}\nosep
  \item If the shuffled series also reaches $R^2 \ge 0.5$, the design fits anything and
        the observed $R^2$ is not read. The verdict is amber, gives the ATT and both
        $R^2$, and sends you to Placebo in space and the other estimators.
  \item Otherwise the verdict is read from the pre-cutoff $R^2$: green at 0.8 or above,
        amber from 0.5, red below; it gives the ATT.
\end{itemize}
\code{mlscm} and \code{bayesian} are too slow to fit twice; their verdict is read from the
pre-cutoff $R^2$ alone. An \code{scm} fit adds a line saying its weights balance the ridge
residuals only.

\screenshot{ui_scm_fit.png}{The augmented SCM (\code{scm}) fit for MY1: ATT $=
-18.14$~\ugm{} over the 11 post-cutoff months. Its pre-cutoff $R^2$ is 0.999, but with
MY1's pre-cutoff values shuffled it still fits 0.9998 (both shown as 1.00), so the verdict
does not read the fit and points to the placebo tests.}{fig:scmfit}

With 53 donors for 27 pre-cutoff months, this design is one where augmented SCM fits
anything. Its top donor carries 99.5\,\% of the weight ($\text{HHI} = 0.99$), but those
weights only balance the ridge model's residuals.

The estimators do not agree on this design. Run on the same panel, cutoff and donors:
\begin{center}\small
\begin{tabular}{@{}lrrrr@{}}
\toprule
Backend & pre-cutoff $R^2$ & shuffled $R^2$ & pre-cutoff RMSE & ATT (\ugm{}) \\
\midrule
\code{scm}    & 0.999  & 1.000   & 0.16  & $-18.14$ \\
\code{abadie} & $-2.94$ & $-3.33$ & 13.20 & $-6.83$ \\
\code{robust} & 0.37   & 0.16    & 5.27  & $-9.48$ \\
\code{did}    & $-0.23$ & $-0.09$ & 7.36  & $-20.67$ \\
\bottomrule
\end{tabular}
\end{center}
The shuffled column is the baseline the \tab{Fit} tab computes: only for \code{scm} does
a series unrelated to the donors fit as well as MY1, so only its pre-cutoff $R^2$ goes
unread. Classic SCM cannot fit MY1 at all: its weights are non-negative and sum to one, so the
synthetic unit cannot exceed the highest donor, and MY1 is above every donor. Robust SCM
fits it only loosely. The four estimates agree on the sign and differ threefold in size.
That is a finding about the design, not a fault of the software: a treated unit outside
its donors' range is the case synthetic control handles worst, and an analysis of it
should report the spread, not one number.

\section{Inference}
\begin{description}
  \item[Run placebo in space] The \tab{Placebo Space} tab draws every donor's placebo
        effect in grey and the treated unit's in red, then the treated effect inside a
        95\,\% band of the placebos. Green at $p \le 0.05$, amber with ``borderline'' up to
        0.15, amber otherwise. It also reports Abadie's RMSPE-ratio test: the treated
        unit's post/pre RMSPE ratio, its rank among the placebos and its p-value.
  \item[Run placebo in time] Fake cutoffs every \emph{Placebo every} steps (7), each with
        at least \emph{Min pre-period} steps (30) before it. A fake post-period must fit
        before the true cutoff, so this needs a long pre-period; the tab says so when
        none fits. Green at $p \le 0.05$; when $p$ is the smallest the number of fake
        cutoffs allows, the verdict says so and suggests a smaller step.
  \item[Run pre-trend + event study] Was the fit parallel before the intervention? A
        slope test on the pre-cutoff residuals --- its slope, p-value and an RMPE drift
        ratio, the fit's own parallel-trends check, which the \tab{Fit} tab reports too
        --- and the effect path on the event-time axis: the pre-period points should
        scatter about zero and the post-period should leave them.
  \item[Run conformal interval] A finite-sample 95\,\% interval for the average
        post-cutoff effect, built from a sub-sampling null distribution of the effect's
        block means (Chernozhukov, W\"uthrich \& Zhu 2021); the number of sub-samples is
        the \emph{Conformal draws} control. Green at $p \le 0.05$.
  \item[Run leave-one-donor-out] The control is refit leaving out each donor in turn:
        how far the remaining weights and the post-cutoff effect move when one donor
        drops. A result carried by a single donor stands out here.
  \item[Uncertainty bands] Jackknife (leave each donor out) or bootstrap ($B$ refits,
        100 by default). The verdict gives the mean post-cutoff effect and on what
        share of post-cutoff dates the band excludes zero; green at half or more.
  \item[Run all units] Fits every unit as if it had been treated, with all the others as
        donors, and ranks them by the size of their post-cutoff effect. The treated unit
        should stand out: green when it is in the top tenth.
\end{description}

\screenshot{ui_scm_placebo.png}{Placebo in space for the ULEZ: no donor's placebo effect
comes near MY1's, so $p = 1/54 = 0.019$, the smallest 53 placebos allow. Every placebo is
fitted by the same estimator, so this comparison is fair even where the pre-period fit
is not informative.}{fig:scmplacebo}

\screenshot{ui_scm_bands.png}{Jackknife bands: the effect is recomputed 53 times, leaving
out one donor each time. The mean post-cutoff effect is $-18.14$~\ugm{}, and the band
excludes zero on 55\,\% of the post-cutoff dates (the verdict says ``days''; the panel is
monthly).}{fig:scmbands}

% ============================ 12. DATA ========================================
\chapter{Getting data in}\label{ch:data}

\section{Input format}\label{sec:input}
\emph{Open Data\ldots} (Ctrl+O) reads CSV and other delimited text
(\code{.csv}, \code{.tsv}, \code{.txt}, \code{.dat}; comma, semicolon, tab, bar or space, a
decimal comma, UTF-8 or a Windows code page --- each told from the file), any of them
gzipped, Excel (\code{.xlsx}, \code{.xlsm}, \code{.xls}), Parquet, Feather and JSON. The table needs:
\begin{itemize}\nosep
  \item its times: a column named \code{date}, \code{datetime}, \code{time},
        \code{timestamp} (any case, or their Chinese words), a date
        column with a time-of-day column beside it, or \code{year}, \code{month},
        \code{day} (and \code{hour}) columns;
  \item the pollutant(s), numeric;
  \item the meteorology and any other predictors, numeric.
\end{itemize}
The times may be written in any of the usual ways: ISO~8601, \code{31/01/2020 13:00}
(day first) or \code{01/31/2020} (month first) --- the order told from the values, and
said under the file's name when they do not tell ---, \code{31.01.2020},
Chinese year--month--day writing, \code{2020013113}, Unix seconds or Excel's serial days. How they
were read is shown under the file's name. One row per timestamp, hourly or coarser. A
table of several sites in one file belongs in the SCM Studio, or in a script with
\code{do\_all\_multisite}. A table without times is refused, with a pointer to the SCM
Studio for panel data. Missing values are allowed: rows without a target are dropped,
missing predictors filled (Section~\ref{sec:prepare}).

\section{The clock: time zone}\label{sec:clock}
The \emph{Time zone} under the file's name is the clock the analysis keeps and its
figures and tables show: the site's local time, or UTC (the default). Times written
without an offset are taken to be on it, as a monitoring network's local-time record is;
times written with one (\code{+08:00}, \code{Z}) are put on it. The hour and the weekday
the model learns from are then the site's own --- a rush hour stays at its hour through
summer time. Meteorology and trajectories fetched on UTC (the Data and Transport Studios)
are put on this clock as they are joined, so an hour of ERA5 meets the hour of the
measurement it belongs to. The choice is kept for the next session; changing it reads the
open file again. In a script, \code{nm.io.read\_time\_series(path, zone=\ldots)} reads a
table the same way.

% The Data Studio's figure is taller than half a page: started at the foot of a
% page it is floated to a float page of its own, so the section opens a fresh one.
\clearpage
\section{Data Studio: air quality networks and meteorology}\label{sec:datastudio}
\emph{Get data} (Ctrl+D) assembles a model-ready table without leaving the app
(Figure~\ref{fig:datastudio}):
\screenshot{ui_data_studio.png}{The Data Studio after a search on the EEA network for
\code{DE}: the stations on the map, three picked and ringed, their rows highlighted in
the table below --- \emph{Fetch \& merge} would download all three into one table.
Network, pollutants, date range and meteorology source on the left; the
\emph{Source settings} group belongs to the non-UK networks.}{fig:datastudio}

\begin{enumerate}\nosep
  \item \textbf{Air quality.} Choose a network. The six UK networks are AURN (the
        national network, about 210 sites), AQE, SAQN, WAQN, NI and LMAM (locally
        managed), about 1500 stations in all; they need neither a country nor a key.
        The \emph{EEA} row searches the European Environment Agency's reporting
        network for one country (ISO-2, e.g.\ \code{DE}): a few hundred to a thousand
        stations, also with no key. \emph{OpenAQ} covers most of the world and wants
        a free API key --- paste it into \emph{Source settings}, or set
        \code{OPENAQ\_API\_KEY}. Tick the pollutants (\PMtwo{}, \NOtwo{} and
        O\textsubscript{3} by default), and \emph{Find stations}. The \tab{Stations}
        tab draws the sites at their coordinates on a map, with the table beneath
        it: a click selects the station under the pointer, a box drawn with the
        mouse selects every station inside it, and Command/Control-click adds
        one. The table follows either way, and the search field narrows both (a
        name, or an official code such as \code{MAN3} for Manchester
        Piccadilly); hovering a point names the station. The wheel zooms about
        the pointer, the $+$ and $-$ buttons step it, the hand button (or the
        middle button) pans, \emph{Fit map} brings the stations back and
        \emph{Deselect} (or Esc) empties the selection; double-click a station
        to zoom to its streets, and right-click for the menu --- select what the
        view holds, copy a station's coordinates. The row above the map counts
        what the search found and how much of it is picked. A basemap lies under the points
        --- OpenStreetMap, Esri's light canvas or satellite imagery, picked on
        the row above the map, key-free and credited below it --- fetched when
        the view is shown.
  \item \textbf{Date range.} The UK archives are whole calendar years, from each
        station's opening; selecting a station snaps the range to its record. The EEA
        record is dated pollutant by pollutant --- its UK stations end in 2019,
        because the UK stopped reporting to the agency after Brexit --- and the range
        snaps to the years the chosen station covers; an OpenAQ location carries its
        last reporting date.
  \item \textbf{Meteorology.} Open-Meteo's ERA5 archive (no key), or the Copernicus CDS
        (needs \code{\textasciitilde/.cdsapirc}). Open-Meteo adds \code{t2m},
        \code{d2m}, \code{rh2m}, \code{sp}, \code{tcc}, \code{tp}, \code{ssrd},
        \code{ws}, \code{wd}, \code{u10} and \code{v10} --- names the \emph{Columns} box
        recognises as meteorology. They come on UTC and are put on the window's clock
        (Section~\ref{sec:clock}) as they are joined.
  \item \textbf{Fetch \& merge} downloads the hourly measurements and the meteorology for
        the site's coordinates and joins them on the hour. With several stations
        selected it does that for each and stacks them into one table, which
        carries a \code{site} column --- a panel the modelling steps can group.
        \tab{Preview} and
        \tab{Table} show the result; \emph{Save merged CSV\ldots} writes it and
        \emph{Send to main window} opens it there.
\end{enumerate}
The UK archives are the openair \code{.RData} files R's \code{importUKAQ()} reads, parsed
without R, so the data match an openair workflow. The EEA path reads the agency's own
station map and downloads one public Parquet file per station and pollutant --- a whole
record in one download. OpenAQ answers per sensor; the \ugm{} sensor is preferred where a
location registers several, so its values share the UK and EEA scale (a ppm-only location
keeps its own unit, in the \code{unit} of its rows), and it carries the six common
pollutants: a ticked NO or NO\textsubscript{x} is skipped for it, with a line in the log.
The Data Studio needs an internet connection; a table sent from it is held in memory, and
the command log says so rather than writing a \code{read\_csv} it could not repeat.


\section{Transport Studio: back-trajectory predictors}\label{sec:transport}
Local meteorology does not say where the air came from. The Transport Studio turns
HYSPLIT back-trajectories into predictors that do, and joins them onto the table
(Figure~\ref{fig:transport}). Its panel has two pages: \emph{Trajectories}, for files
already made, and \emph{HYSPLIT}, for running them.

\paragraph{Trajectory files.} \emph{Add files\ldots} takes existing HYSPLIT
\code{tdump} files (no HYSPLIT installation is needed to read them). \emph{Site name}
tags the result with a \code{site} column, so batches for several receptors, built one
at a time, each fill only their own site's rows. \emph{Null trajectories shorter than}
(72~h) blanks trajectories that stopped early because the meteorology ran out or the air
left the domain, which would otherwise pass for short-range transport.

\paragraph{Source regions.} Optional regions, as bounding boxes (name, lon/lat min and
max) or polygons from GeoJSON or Shapefile files (\code{geo} extra). Each adds a
residence-time feature: the fraction of the trajectory spent inside. Overlapping regions
each count an endpoint inside both.

\paragraph{Features.} Every trajectory becomes one row of: \code{traj\_n\_endpoints} and
\code{traj\_age\_max\_h} (quality: how far back it reached), \code{traj\_dist\_km}
(straight-line reach), \code{traj\_pathlen\_km}, \code{traj\_speed\_kmh},
\code{traj\_inflow\_deg} (bearing from the receptor to the origin),
\code{traj\_height\_mean} and \code{traj\_height\_min}, one
\code{traj\_resid\_}\emph{region} per region, and --- where the \code{tdump} carries
them --- along-path rainfall, mixing depth, humidity, pressure and temperature.

\paragraph{HYSPLIT (advanced).} With a local \code{hyts\_std}, the studio can run the
trajectories itself, on its \emph{HYSPLIT} page, for a table of receptors, at a chosen
height, hours back and time step, with GDAS1 meteorology it downloads or ARL files already
on disk. The features are built with the source regions set on the \emph{Trajectories}
page.

\paragraph{Joining.} \emph{Send to main window} joins the features onto the open table:
each row takes the latest trajectory that has \emph{already arrived}, within one release
interval --- a nearest-time join would give 04:00 the 06:00 trajectory, air that has not
yet arrived --- and the new \code{traj\_*} columns are ticked as predictors. The
trajectories' times, UTC as HYSPLIT writes them, are put on the window's clock first
(Section~\ref{sec:clock}); the studio is given the table's times in UTC to start them from.

\screenshot{ui_transport.png}{The two days of \code{tdump} files that ship with \nm{},
built for receptor MY1: 9 features for 8 receptor times, released every six hours.
Inflow turns from the east (80$^\circ$) to the west (270$^\circ$) over 1--2 January
2020.}{fig:transport}

\caution{Trajectories released every six hours are constant over six hourly rows. Under a
random split a tree model uses them as a time stamp and its test $R^2$ is inflated
(Section~\ref{sec:split}). Judge a transport-aware model with \code{month\_ts}. To split
\PMtwo{} into local-weather and transport contributions, give the meteorological
decomposition two groups, \code{local} and \code{transport}
(Section~\ref{sec:metdecomp}).}

\section{Polar plot: where a pollutant comes from}\label{sec:polar}
\tab{Analysis $\rightarrow$ Polar Plot\ldots} draws a variable of the open table by wind
direction and speed, as openair's \code{polarPlot} does (Figure~\ref{fig:polar}): its
mean, median, 95th percentile or maximum in bins of 10$^\circ$ of direction and of speed,
on a compass, north at the top. The wind columns are found by name (\code{ws} and
\code{wd}), else made from \code{u10} and \code{v10}, and can be chosen. In a script:
\code{nm.polar\_plot(df, target="NO2", ws\_col="ws", wd\_col="wd")}.

\screenshot{ui_polar.png}{Mean \NOtwo{} at MY1 by wind direction and speed: highest with
winds from the south and south-west, lowest with winds from the north and
north-east.}{fig:polar}

% ============================ 13. EXPORT ======================================
\chapter{Exporting and reproducing}\label{ch:export}

\section{Tables and reports}
\begin{description}
  \item[Export result CSV] (Ctrl+E) The table behind the current tab: the loaded data on
        \tab{Data}, the statistics on \tab{Model}, the normalised series with its
        quantiles, the decomposition, every rolling window, the multi-scale series, the
        partial-dependence values.
  \item[HTML report] (File $\rightarrow$ Export HTML Report, or the button; pick the
        report's language from its menu.) The current tab's result as one self-contained HTML file --- figures
        embedded, no scripts or external links --- with a summary of the model and the
        result, and its provenance: \nm{}'s version, the time, the platform and Python
        version, the host and user names, the shape and hash of the data, the panel's
        settings and their hash, and the seed. Remove the host and user names before
        sharing a report if they should not travel with it.
\end{description}

\section{The command log}\label{sec:commandlog}
Every step that decides a result --- opening a table, training, normalising,
decomposing, a rolling run, a multi-scale run, partial dependence --- is written to
\tab{View $\rightarrow$ Command Log\ldots} as the \nm{} call that does the same, with the
arguments the window actually used. An entry is written only once its task has
succeeded, from the call that was made, not from what the panel shows now.
\emph{Copy} puts the log on the clipboard, \emph{Save as Script\ldots} (also \tab{File
$\rightarrow$ Save Session as Script\ldots}) writes it as a Python file that repeats the
session, and \emph{Clear} empties it.

\screenshot{ui_command_log.png}{The Command Log after this manual's session, its last
entries in view: the multi-scale run and the partial dependence, as the calls that
repeat them.}{fig:commandlog}

The session behind this manual's figures saves as (paths shortened, imports and the
multi-scale and partial-dependence calls left out):
\begin{lstlisting}[language=Python]
# 21:09:36  Load example: London Marylebone Road NO2 (2020)
df = nm.datasets.load_my1()

# 21:10:09  Train model
df_prep, model = build_model(df, 'NO2', backend='flaml',
                             covariates=['ws', 'wd', 'temp', 'RH', 'atmos_pres', 'blh', 'tcc', 'tp', 'ssrd', 'date_unix', 'day_julian', 'weekday', 'hour'],
                             split_method='random', train_fraction=0.75,
                             model_config={'time_budget': 30, 'estimator_list': ['lgbm']},
                             seed=7654321)

# 21:10:11  Normalise
normalised = normalise(df_prep, model,
                       covariates=[...the same 13...],
                       variables_resample=['ws', 'wd', 'temp', 'RH', 'atmos_pres', 'blh', 'tcc', 'tp', 'ssrd'],
                       n_samples=300, n_cores=None,
                       return_quantiles=(0.05, 0.25, 0.75, 0.95))

# 21:10:22  Decompose
decomposition = decompose(df_prep, model, method='meteorology',
                          covariates=[...the same 13...], n_samples=100,
                          groups={'Wind': ['ws', 'wd'],
                                  'Temperature and humidity': ['temp', 'RH'],
                                  'Boundary layer and pressure': ['atmos_pres', 'blh'],
                                  'Sky and rain': ['tcc', 'tp', 'ssrd']})

# 21:10:25  Rolling normalisation
rolled = rolling(df_prep, model, covariates=[...],
                 variables_resample=[...the nine...],
                 window_days=14, rolling_every=7, n_samples=50)
\end{lstlisting}
Run again, it gives the same normalisation and decompositions from the same model; the
model itself can differ slightly, since FLAML's search is bounded by wall-clock time
(Section~\ref{sec:train}).

\section{Figures for publication}\label{sec:print}
A figure on screen is drawn at the tab's size; shrunk into a journal's 89~mm column its
lettering shrinks with it. \emph{For print\ldots}, on every figure's toolbar, lays the
figure out again at the size the journal prints it, each text at the size its guide gives
that kind of text --- a title, an axis label, a tick label, a legend --- with the margins
and the ticks worked out anew for that size, not scaled to it (Figure~\ref{fig:print}):
\begin{description}\nosep
  \item[Journal] Elsevier, Nature, Copernicus (ACP, AMT, \ldots) or ACS (ES\&T), each with
        what its own artwork guide asks, and a link to that guide.
  \item[Width] The journal's columns (Nature 89 or 183~mm; Elsevier 90, 140 or 190~mm;
        ACS 3.33 or 7~in; Copernicus at least 80~mm), or a width of your own. A figure of
        one column of panels starts at the narrowest, panels side by side at the widest.
  \item[Height] As the figure needs --- its own proportion, raised until every panel has
        at least 12~mm for its data, within the journal's greatest height, and lowered to
        what is drawn where panels of a shape of their own (roses side by side) would
        leave the rest blank --- or set.
  \item[Panel labels] None, a~b~c, (a)~(b)~(c), A~B~C or (A)~(B)~(C), bold or not, from any
        letter (a figure that is panels c and d of a composite starts at c); each column of
        panels has its letters in a line.
  \item[Line width, Marker size] The data lines' stroke; a time series takes round
        markers when the size is raised. Reference lines keep theirs.
  \item[Legend] Where it sits (or drag it), its text size, its columns, its frame. A legend
        laid out in a row wider than the printed figure is given fewer columns.
  \item[Formulas] Written as chemistry writes them: NO$_2$, PM$_{2.5}$, SO$_4^{2-}$.
  \item[Titles] The figure's and the panels' titles are left to the caption, as journals
        ask; ticked, they are drawn.
  \item[Raster resolution] 300, 600 or 1200~dpi, for TIFF and PNG.
\end{description}
Under the settings, the checks read the figure as it is laid out against the journal's
guide: its width and height, its smallest text and its smallest sub- or superscript
(Nature asks for 5~pt at least and 7~pt at most, Elsevier 7~pt, ACS 8~pt), the face it is
drawn in (Arial or Helvetica where the guide names them), tick labels that run into one
another, and how many panels it has (ACS: four at most); a check the figure fails is shown
in red, and all are read again as the settings change. The preview fits the window, or
shows the figure at its actual printed size. On the preview, double-click a title, an axis
label or a panel letter to edit its text and size, a tick's numbers to size and turn them
(and give that tick a text of its own), the legend to place it, or a colour bar to put it
on another side of its panel and give its width and height on the page in millimetres;
drag the legend or a colour bar where it reads best.
\emph{Save\ldots} writes the formats the journal takes --- PDF, EPS or SVG with the text
kept as editable text and the fonts embedded, TIFF or PNG at the resolution chosen --- at
exactly the size shown, on a white page. While the dialog is open the figure is in the
report's language (the one last picked under File $\rightarrow$ Export HTML Report). When it closes the figure goes back
to its tab as it was; the journal, width, panel labels, formulas and resolution are kept
for the next figure. Check the journal's current guide before submitting.

\screenshot{ui_print.png}{The normalised series laid out for a Nature double column:
183~mm wide, its panels lettered, beside what Nature's guide asks of it --- a subscript
at 4.9~pt, below the 5~pt Nature asks, is the one check it fails: Nature's 7~pt ceiling
leaves no room for more.}{fig:print}

\section{Saving and reopening a session}\label{sec:saved}
\tab{File $\rightarrow$ Save Results\ldots} (Ctrl+S) writes the session to one file
(\code{.normet}): the table, the panel's settings, the model, every step's result and the
command log. \tab{Load Results\ldots} reads one back into the window as it was --- each
result is handed to its step's own handler, so its tab, verdict and figure are drawn again
--- and a notice says what was restored. After each step that finishes, the window also
keeps the session as \code{\textasciitilde/.normet/last\_session.normet};
\tab{Reopen Last Session} reads it, should the window have been closed, or have crashed,
before the results were saved. A Chronos-2 checkpoint is not written (it is the
pretrained model, about 500~MB): load it again to go on. A session file is a Python pickle:
open only one you or the window wrote.

\tab{File $\rightarrow$ Load Session Script\ldots} runs a script the command log wrote
(Section~\ref{sec:commandlog}) and brings its results into the window the same way: its
table, model and results, and the panel set from the calls it made. The script is Python
and runs as it is written, so the window asks first.

\section{A Methods paragraph}\label{sec:methods}
\tab{File $\rightarrow$ Methods Paragraph\ldots} writes what the session ran as a paper's
Methods says it, ready to copy (Figure~\ref{fig:methods}): the record (its length, span and
step), the model (backend, search, predictors, split and seed) and its held-out fit, then
each step run --- the normalisation's draws and resampled variables, the decomposition, the
rolling windows, the multi-scale widths, the partial-dependence variables --- with the
references the methods rest on. It is written from each step's call as it ran, not from
the panel as it stands, and in the window's language. Check it against what you did.

\screenshot{ui_methods.png}{The Methods paragraph of this manual's session: the record,
the FLAML model and its held-out fit, and each step that was run, with its
references.}{fig:methods}

\section{Configurations and history}\label{sec:config}
\tab{File $\rightarrow$ Save Config\ldots} writes every setting of the panel to JSON ---
target, predictors, backend, split, budget, estimators, seed, and every step's settings
--- and \tab{Load Config\ldots} sets the panel from one; the data are not part of it.
\tab{File $\rightarrow$ Run History\ldots} lists every run with its time, task, target, backend, rows and a
detail (the held-out $R^2$ for a training run, the samples for a normalisation).

% ============================ 14. ASSISTANT ===================================
\chapter{The Assistant}\label{ch:assistant}

The Assistant is a language model working in the window. It is open on the right when
\nm{} starts; \tab{View $\rightarrow$ Assistant} (Ctrl+Shift+A) closes it and opens it
again. It reads what the window holds --- the table, the model's held-out fit, the
normalised series, the decomposition --- as \nm{}'s own readings of them, in text, and
runs the window's steps the way their buttons do, so what it does appears on screen as
your own runs do and goes into the command log.

The conversation reads as a chat does: your messages on the right, the replies as their
Markdown reads (code shaded, tables ruled), each tool the Assistant calls one quiet line
--- its dot grey while it runs, green when it answered, red when it was refused, the
refusal under it --- whose result opens with a click. Ctrl/Cmd $+$ and $-$ (or the wheel
with the key held) size its text, and Ctrl/Cmd 0 resets it; the size is kept. You write
in the rounded box at its foot, which grows with what you write: Enter sends, Shift+Enter
starts a new line. \emph{Send}, \emph{Stop} and \emph{New chat}, which starts the
conversation again, are inside the box.

Where it helps most is what to model and how to read the result. It \emph{proposes} a
target and predictors as a card with an \emph{Apply} button, saying what it left out and
why; nothing in the panel changes until you click Apply, and it trains or normalises only
once you have agreed. It is told to quote only numbers the tools return.

\section{Full automatic}\label{sec:autoassistant}
\emph{Auto}, beside \emph{Key\ldots}, lets the Assistant work without waiting for you.
With it on, the target and predictors it proposes are applied at once (the
card says \emph{Applied}), and it works the analysis through in the window's order ---
the data, the set-up, training, normalisation, decomposition --- reading each result as it
comes. A model that stops with steps left is told what is left and goes on, up to eight
times; two reminders in a row that change nothing end the turn, as a model that cannot do
a step should not circle. A call to the model that fails, or an empty reply, is made again,
twice at most. Each reminder and retry is a line in the conversation. \emph{Stop} ends it
at any time; the choice is kept for the next session, and \emph{What full automatic
does\ldots} under the button's arrow says all this. A command-line agent has the set-up
applied at once too, and works the chain as far as its own loop allows.

\section{Choosing a model}
\begin{description}
  \item[An API model] named \code{provider:model}: \code{anthropic}, \code{openai},
        \code{deepseek}, \code{gemini}, or a local \code{ollama} or \code{lmstudio}
        server --- for example \code{anthropic:claude-sonnet-5-5},
        \code{deepseek:deepseek-chat}, \code{ollama:qwen3:8b}. Its key is read from the
        macOS keychain (\emph{Key\ldots} stores it there) or the provider's environment
        variable (\code{ANTHROPIC\_API\_KEY}, \code{OPENAI\_API\_KEY},
        \code{DEEPSEEK\_API\_KEY}, \code{GEMINI\_API\_KEY}), and written nowhere else.
        Local servers need no key.
  \item[A command-line agent] signed in to your own plan: Claude Code, Codex, Google
        Antigravity's \code{agy} or GitHub Copilot, with the effort it thinks at. The
        window serves its tools to the agent over MCP on 127.0.0.1 with a token drawn for
        that window, and gives it no shell and no files of its own.
\end{description}
What the Assistant sends to a remote model is the readings --- summaries, statistics,
captions --- and your messages, not the table.

% ============================ 15. SCRIPTING ===================================
\chapter{Scripting: API, CLI and MCP}\label{ch:script}

\section{Python API}
Everything the window does is a function of the \code{normet} package. The one-call
pipeline prepares, trains and normalises:
\begin{lstlisting}[language=Python]
import normet as nm

df = nm.datasets.load_my1()
met = ["ws", "wd", "temp", "RH", "atmos_pres", "blh", "tcc", "tp", "ssrd"]
predictors = met + ["date_unix", "day_julian", "weekday", "hour"]

out, model, df_prep = nm.do_all(
    df=df, target="NO2", backend="flaml",
    covariates=predictors, variables_resample=met,
    n_samples=300, n_cores=4,
)
print(nm.modStats(df_prep, model))   # training, testing and all rows
\end{lstlisting}
Step by step:
\begin{lstlisting}[language=Python]
df_prep, model = nm.build_model(df, "NO2", backend="flaml", covariates=predictors,
                                model_config={"time_budget": 120})
norm = nm.normalise(df_prep, model, covariates=predictors,
                    variables_resample=met, n_samples=300,
                    return_quantiles=[0.025, 0.975])
auto = nm.normalise_auto(df_prep, model, covariates=predictors,
                         variables_resample=met)       # stops when converged
emi = nm.decompose(df_prep, model, covariates=predictors, method="emission")
mets = nm.decompose(df_prep, model, covariates=predictors, method="meteorology",
                    groups={"wind": ["ws", "wd"], "rest": met[2:]})
roll = nm.rolling(df_prep, model, covariates=predictors, variables_resample=met)
curves = nm.pdp(df_prep, model, var_list=["ws", "temp"])
\end{lstlisting}
Further functions, all exported at the top level:
\begin{center}\small
\begin{tabular}{@{}p{7.4cm}p{8.0cm}@{}}
\toprule
Function & Does \\
\midrule
\code{do\_all\_unc}       & the pipeline over \code{n\_models} seeds, with confidence bands \\
\code{do\_all\_multisite} & per-site pipelines in parallel \\
\code{cv\_score}          & walk-forward cross-validation \\
\code{run\_scm}, \code{scm\_all} & a synthetic control; every unit as if treated \\
\code{placebo\_in\_space}, \code{placebo\_in\_time} & placebo tests \\
\code{uncertainty\_bands}, \code{effect\_bands\_space} & jackknife/bootstrap and placebo bands \\
\code{scm\_diagnostics}, \code{conformal\_effect\_interval} & fit diagnostics; conformal interval \\
\code{bayesian\_scm}      & Bayesian SCM (\code{bayes} extra) \\
\code{detect\_events}, \code{anomaly\_scores} & events in a normalised series \\
\code{normalise\_plot}, \code{decomposition\_stack}, \code{scm\_dashboard} & the window's figures \\
\code{make\_run}, \code{save\_run}, \code{generate\_html\_report} & provenance and reports \\
\code{trend}, \code{mann\_kendall}, \code{theil\_sen} & trend analyses with the block-bootstrap correction (the \tab{Trend} tab, Section~\ref{sec:trend}) \\
\code{detect\_trend\_breaks} & segmented trend: where the level or slope breaks (BIC + block-bootstrap gate) \\
\code{decompose\_uncertainty} & the total spread attributed to weather, model form and their interaction (the \tab{Uncertainty} tab, Section~\ref{sec:uncertainty}) \\
\code{make\_memory}       & on-disk caching \\
\code{Chronos2Estimator}  & zero-shot de-weathering (\code{foundation} extra) \\
\bottomrule
\end{tabular}
\end{center}
The configuration of each long function is also a dataclass (\code{NormaliseConfig},
\code{DecomposeConfig}, \code{RollingConfig}); keywords override its fields. The guides
at \url{https://normet.readthedocs.io} cover each area with worked examples, and the
notebooks in the repository's \code{notebooks/} folder reproduce the paper's case
studies.

\section{Command line}\label{sec:cli}
With the \code{cli} extra, \code{normet} runs the pipelines on a CSV or Parquet file:
\begin{lstlisting}[language=bash]
normet do-all my1.csv --target NO2 \
    --covariates ws,wd,temp,RH,atmos_pres,blh,tcc,tp,ssrd,date_unix,day_julian,weekday,hour \
    --resample-vars ws,wd,temp,RH,atmos_pres,blh,tcc,tp,ssrd \
    --n-samples 300 --out normalised.csv --run-path my1_run
normet report my1_run.joblib --out my1_report.html
normet decompose my1.csv --target NO2 --covariates ... --method emission --out emi.csv
normet scm ulez_panel.csv --treated MY1 --cutoff 2019-04-01 --date-col date \
    --unit-col code --outcome-col NO2_dw --backend scm --out scm.csv
normet cv my1.csv --target NO2 --covariates ... --n-splits 5 --out cv.csv
normet deweather site.csv --target PM2.5 --met-vars t2m,blh,u10,v10 --out dw.csv
normet info                     # version and the backends installed
\end{lstlisting}
Every command takes \code{--config file.yaml}, whose keys are the options' names; an
option on the command line overrides the file. \code{--run-path} also saves a
\code{NormetRun} archive (\code{.joblib} with a \code{.meta.json} sidecar) that
\code{normet report} turns into HTML or Markdown. The command line's \code{scm} offers
\code{scm}, \code{mlscm}, \code{abadie}, \code{did} and \code{mcnnm}.

\section{AI assistants over MCP}\label{sec:mcp}
With the \code{mcp} extra, \code{normet-mcp} serves \nm{} to an AI assistant over the
Model Context Protocol:
\begin{lstlisting}[language=bash]
claude mcp add normet -- normet-mcp          # Claude Code
# Claude Desktop, claude_desktop_config.json:
#   {"mcpServers": {"normet": {"command": "normet-mcp"}}}
\end{lstlisting}
The assistant can list the examples, inspect a table, start \code{do-all},
\code{decompose}, \code{scm} and \code{cv} jobs, follow and cancel them, read a run or a
result, draw a figure and write a report. Each fit is a \nm{} command run in its own
process and folder (\code{\$NORMET\_MCP\_HOME}, \code{\textasciitilde/.normet/mcp} by
default), so its command line is its record and runs again without the assistant. The
assistant is given \nm{}'s readings of a table or result, not the data; files are read
only under \code{\$NORMET\_MCP\_ROOTS} (your home folder by default). The server asks the
assistant to propose the target, predictors and donor pool with its reasons, to let you
confirm them, and to quote only the numbers the tools return.

% ============================ APPENDICES ======================================
\appendix

\chapter{Defaults}\label{app:defaults}
\begin{center}\small
\begin{longtable}{@{}p{4.9cm}p{4.3cm}p{5.8cm}@{}}
\toprule
Setting & Window & Library \\
\midrule
\endhead
Split method / training fraction & random / 0.75 & random / 0.75 \\
FLAML time budget              & 60 s & 90 s \\
FLAML estimators               & lgbm & lgbm \\
FLAML metric                   & $R^2$ & $R^2$ \\
LightGBM search trials         & 20 & --- \\
Seed                           & 7654321 & 7654321 \\
Normalisation samples          & 300 (Chronos-2: 8) & 300 (Chronos-2: 8) \\
Resampled variables            & recognised meteorology & every non-calendar predictor \\
Quantiles                      & 5, 25, 75, 95\,\% when ticked & none \\
Decomposition samples          & 300 & 300 \\
Attribution                    & auto & sequential (Shapley with groups) \\
Rolling window / step / samples & 14 d / 7 d / 100 & 14 d / 7 d / 300 \\
Multi-scale widths             & 14 / 90 / 365 d & --- \\
PDP grid points / range        & 50 / 1st--99th percentile & 50 / 1st--99th percentile \\
\code{normalise\_auto}          & --- & batches of 100, P95 RSE $<$ 3\,\%, 3 in a row, $\le$ 5000 \\
Chronos-2 context / horizon    & 2048 / 168 rows & 2048 / 168 rows \\
SCM backend                    & scm & scm \\
Placebo in time: min pre / every & 30 / 7 steps & --- \\
Uncertainty bands              & jackknife, $B = 100$ & jackknife, $B = 200$, 80\,\% donors \\
\bottomrule
\end{longtable}
\end{center}

\chapter{Where normet keeps its files}\label{app:files}
\begin{center}\small
\begin{tabular}{@{}ll@{}}
\toprule
What & Where \\
\midrule
Recent tables (eight)        & \code{\textasciitilde/.normet/gui\_recent\_files.json} \\
Run history                  & \code{\textasciitilde/.normet/gui\_run\_history.jsonl} \\
Window geometry, appearance, language, Assistant model & the system's settings store (\code{apai-sys/normet-gui}) \\
Font cache (installers)      & \code{\textasciitilde/.normet/matplotlib} \\
The last session (after each step) & \code{\textasciitilde/.normet/last\_session.normet} \\
Assistant API keys           & macOS keychain, service \code{normet} \\
MCP jobs                     & \code{\$NORMET\_MCP\_HOME}, default \code{\textasciitilde/.normet/mcp} \\
Chronos-2 weights            & the Hugging Face cache \\
\bottomrule
\end{tabular}
\end{center}
Nothing else is written unless you save it. The print dialog's choices, the Assistant's
text size and the report's language are kept in the settings store with the rest.

\chapter{Troubleshooting}\label{app:trouble}
\begin{description}
  \item[``The CSV must contain a 'date' column''] Rename the timestamp column to
        \code{date}. Panel data (several units) belongs in the SCM Studio.
  \item[The model's verdict is red] The predictors do not explain the target. Check that
        the meteorology is ticked and on the same time base as the target, lengthen the
        budget, or add estimators. Normalising with a poor model returns little more than
        the observed series.
  \item[A meteorological column is not ticked] Only the names in
        Section~\ref{sec:columns} are recognised; tick others by hand, on the \tab{Data}
        and \tab{Normalise} pages.
  \item[Exact Shapley values over $k$ features need $2^k$ normalisations] More than ten
        variables by Shapley: give groups, or use sequential attribution.
  \item[``\ldots{} not among the variables trained on''] A group (from a configuration
        read in) names a variable the model was not trained on; the message names it.
  \item[``'rdata' is required but not installed''] \emph{Find stations} reads the UK
        archives with \code{rdata}, part of the \code{data} extra:
        \code{pip install "normet[data]"}. The installers include it.
  \item[Multi-scale asks for the normalisation first] It differences against the
        normalised series.
  \item[A multi-scale band was skipped] Its width leaves fewer than two windows on the
        record; lower it.
  \item[Rolling and PDP are disabled] They need an AutoML model; the Chronos-2 backend
        has no response surface to vary.
  \item[Chronos-2 is greyed out] Install \code{normet[foundation]} (and
        \code{normet[finetune]} to fine-tune).
  \item[Chronos-2's verdict is red] Shuffling the weather barely moves its forecast: it is
        autoregressing. Check the meteorological columns, or use an AutoML backend.
  \item[SCM: ``The pre-treatment fit is no check on this design''] A series shuffled out
        of any relation to the donors fits the pre-period as well as the treated unit
        does: there are too many donors for the pre-cutoff dates. Judge the effect by
        Placebo in space and by other estimators, or use fewer, better-chosen donors (the
        similarity ranking can help) or a longer pre-period.
  \item[SCM: ``Donor pool too small''] Tick at least two donors.
  \item[SCM: no valid fake cutoffs] The post-period must fit before the true cutoff with
        \emph{Min pre-period} steps before it; lower \emph{Min pre-period} or
        \emph{Placebo every}, or move the cutoff later.
  \item[A task runs too long] \emph{Cancel} frees the window at once; the computation
        finishes in the background and its result is discarded. Lower the samples, the
        budget, or the number of variables decomposed.
  \item[A notice asks for a step first] It names what is missing and carries the step on
        a button; it goes by itself once the step has been taken.
  \item[The window closed before the results were saved] \tab{File $\rightarrow$ Reopen
        Last Session} brings back the session as it was after its last step.
  \item[The parameter panel is missing] It is folded away: \emph{Expand Panel} on the
        action bar, or Ctrl+B.
  \item[A task failed] The dialog shows its last line; the full traceback is in the
        \tab{Log} dock.
  \item[macOS: ``Normet cannot be opened''] The app is unsigned: right-click it and choose
        \emph{Open} once (Section~\ref{ch:install}).
\end{description}

\end{document}
